% ============================================================================
% APPENDIX XVIII: LIT-FEHR-METHOD - Integration vs. Falsifiability
% ============================================================================
% Category: LIT (Literature/Critical Analysis)
% Version: 1.0
% Status: Complete
% Language: English
% Prerequisites: K (LIT-FEHR), XVI (LIT-CRITIQUE), S (PREDICT-FALSIFIABLE)
% ============================================================================

\begin{chapter}{Fehr's Methodological Challenge: Integration vs. Falsifiability}
\label{app:lit-fehr-method}

% ============================================================================
% HEADER BLOCK
% ============================================================================
\begin{tcolorbox}[colback=gray!5!white,colframe=gray!75!black]
\small
\textbf{Category:} LIT (Literature/Critical Analysis) \\
\textbf{Code:} XVIII \\
\textbf{Version:} 1.0 \\
\textbf{Status:} Complete \\
\textbf{Dependencies:} K (LIT-FEHR), XVI (LIT-CRITIQUE), S (PREDICT-FALSIFIABLE), T (Metatheory) \\
\textbf{Scientific Purpose:} Formalization of Fehr's methodological critique and the EBF response
\end{tcolorbox}

% ============================================================================
% CROSS-REFERENCE MAP
% ============================================================================
\begin{tcolorbox}[colback=blue!5!white,colframe=blue!75!black,title=\textbf{Cross-Reference Map}]
\small

\textbf{CORE Appendices (10C Framework):}
\begin{itemize}[noitemsep]
    \item \textbf{AAA (CORE-WHO):} Welfare hierarchy --- $\gamma$ applies at all levels $L$
    \item \textbf{B (CORE-HOW):} Complementarity structure --- $\gamma$-matrix definition
    \item \textbf{C (CORE-WHAT):} Utility dimensions --- FEPSDE interactions via $\gamma$
    \item \textbf{V (CORE-WHEN):} Context --- $\Psi$ modulates $\gamma(\Psi)$
    \item \textbf{BBB (CORE-WHERE):} Parameters --- $\gamma$ calibration values
    \item \textbf{AU (CORE-AWARE):} Awareness --- awareness of complementarities
    \item \textbf{AV (CORE-READY):} Willingness --- $\gamma$ in WAX formula
    \item \textbf{AW (CORE-STAGE):} BCJ --- stage-specific $\gamma$ patterns
\end{itemize}

\textbf{FORMAL Appendices:}
\begin{itemize}[noitemsep]
    \item \textbf{A (FORMAL-AXIOMS):} Axiom system --- mathematical foundation for $\gamma$
    \item \textbf{D (FORMAL-CSTAR):} Reference structure --- $C^*$ proofs, coherence
\end{itemize}

\textbf{LIT/PREDICT Appendices:}
\begin{itemize}[noitemsep]
    \item \textbf{K (LIT-FEHR):} Fehr's research program and empirical contributions
    \item \textbf{XVI (LIT-CRITIQUE):} Ten canonical arguments against complementarity
    \item \textbf{S (PREDICT-FALSIFIABLE):} Ten testable predictions from EBF
    \item \textbf{T (Metatheory):} Framework philosophy and scientific status
\end{itemize}

\textbf{METHOD Appendices:}
\begin{itemize}[noitemsep]
    \item \textbf{VC (METHOD-VALIDATE):} Data consistency validation --- validates formula parameters (DCV-3, DCV-4)
\end{itemize}

\textbf{Key Insight:} This appendix provides the \textit{methodological foundation} that disciplines all $\gamma$-claims across the entire 10C framework.
\end{tcolorbox}

% ============================================================================
% CHAPTER LINKAGE
% ============================================================================
\begin{tcolorbox}[colback=purple!5!white, colframe=purple!75!black,
    title=Chapter Linkage]

\textbf{Primary Chapters:}
\begin{itemize}[nosep]
\item Chapter 5: Complementarity $\leftrightarrow$ Methodological Note Box
\item Chapter 9: Context $\leftrightarrow$ Local validity of $\gamma$ claims
\end{itemize}

\textbf{Secondary Chapters:}
\begin{itemize}[nosep]
\item Chapter 1: Introduction --- Positions EBF as instrument, not theory
\item Chapter 3: Utility Theory --- Arrow-Debreu as theoretical scaffold
\end{itemize}

\textbf{Reading Guide:}
\begin{itemize}[nosep]
\item For conceptual overview: Read Sections 1--4 (Taxonomy, Exclusion Principle)
\item For philosophy of science: Read Sections 9--10 (Three Levels, Three Traditions)
\item For practical application: Read Sections 17--19 (Positioning, Balance, Practice)
\item For collaboration guidance: Read Section 11 (Practical Collaboration Guide)
\end{itemize}

\end{tcolorbox}

% ============================================================================
% ABSTRACT
% ============================================================================
\begin{abstract}
This appendix formalizes a meta-level critique of complementarity frameworks that is implicit in Ernst Fehr's methodological position but rarely stated explicitly: \textbf{the tension between model integration and scientific falsifiability}. Unlike the ten canonical arguments (Appendix CRT), which concern measurement, identification, or complexity, Fehr's deeper concern is epistemological: if a framework can integrate all models, it loses the disciplinary power that makes models scientific. This appendix (1) reconstructs Fehr's position from his methodological commitments, (2) shows how complementarity ($\gamma$) can either confirm or refute this concern depending on its use, and (3) proposes a resolution: $\gamma$ as an \textit{exclusion principle}, not an integration principle. The central claim: complementarity is scientifically legitimate precisely when it tells us what is \textit{not} complementary.
\end{abstract}

% ============================================================================
% FUNDAMENTAL QUESTION
% ============================================================================
\begin{tcolorbox}[colback=red!5!white,colframe=red!75!black,title=\textbf{The Fundamental Question}]
\begin{center}
\large
\textbf{If complementarity can integrate all models, does it explain anything?}
\end{center}

\vspace{0.5em}
This is not a question about measurement or identification (Appendix CRT), but about the \textit{scientific status} of integrative frameworks themselves.
\end{tcolorbox}

% ============================================================================
% QUICK REFERENCE
% ============================================================================
\begin{tcolorbox}[colback=yellow!5!white,colframe=yellow!75!black,title=\textbf{Quick Reference: The Fehr Position}]
\small
\textbf{Core Thesis:}
\begin{quote}
A model is scientific if and only if it can fail. Model integration threatens falsifiability because any deviation can be attributed to ``missing complementarity.''
\end{quote}

\textbf{Key Distinction:}
\begin{itemize}[noitemsep]
    \item \textbf{Integration Principle:} $\gamma$ explains why everything fits together
    \item \textbf{Exclusion Principle:} $\gamma$ specifies where complementarity does \textit{not} hold
\end{itemize}

\textbf{Resolution:} EBF is scientifically valid when it commits to $\gamma = 0$ null hypotheses.
\end{tcolorbox}


% ============================================================================
% PART I: RECONSTRUCTING FEHR'S POSITION
% ============================================================================
% -----------------------------------------------------------------------------
% APPENDIX SCOPE BOX
% -----------------------------------------------------------------------------

\begin{tcolorbox}[
    colback=orange!5!white,
    colframe=orange!75!black,
    title={\textbf{Appendix Scope: Ziel / In-Scope / Out-of-Scope / Constraints}},
    fonttitle=\bfseries
]

\textbf{Ziel (Objective):}
\begin{itemize}[nosep]
    \item If complementarity can integrate all models, does it explain anything?
\end{itemize}

\textbf{In-Scope:}
\begin{itemize}[nosep]
    \item Reconstructing Fehr's Methodological Position
    \item Why Complementarity Maximizes the Integration Problem
    \item The Resolution: $\gamma$ as Exclusion Principle
    \item Implications for EBF Methodology
    \item The Deeper Philosophical Issue
\end{itemize}

\textbf{Out-of-Scope (delegiert an andere Appendices):}
\begin{itemize}[nosep]
    \item PREDICT-FALSIFIABLE $\rightarrow$ Appendix SUN
    \item WHO $\rightarrow$ Appendix WHO
    \item HOW $\rightarrow$ Appendix HOW
    \item WHAT $\rightarrow$ Appendix WAT
    \item WHEN $\rightarrow$ Appendix CTW
\end{itemize}

\textbf{Constraints (Anwendungsgrenzen):}
\begin{itemize}[nosep]
    \item \textbf{Does NOT apply:} Pure measurement questions (use Appendix AN), single-factor models (no aggregation needed), non-behavioral domains (physics, chemistry)
    \item \textbf{Prerequisites:} Understanding of EXC-1 to EXC-6 axioms, familiarity with potential outcomes notation, basic knowledge of causal identification (RCT/IV/DiD)
    \item \textbf{Limitations:} Does not resolve which $\gamma_{ab} = 0$ commitments EBF should make (Open Issue 1); causal identification requires data (theoretical models use T1, not E1)
\end{itemize}

\textbf{Lieferobjekte (Deliverables):}
\begin{itemize}[nosep]
    \item 11 Table(s)
    \item 2 Equation(s)
    \item Worked Example(s)
\end{itemize}

\end{tcolorbox}


\section{Reconstructing Fehr's Methodological Position}
\label{sec:fehr-position}

\subsection{The Explicit and Implicit Fehr}

Ernst Fehr's published work (Appendix FEH) demonstrates rigorous experimental methodology, but rarely engages in explicit philosophy of science. However, his methodological commitments imply a coherent epistemological position that we can reconstruct.

\subsubsection{What Fehr Explicitly Endorses}

\begin{enumerate}
    \item \textbf{Experimental rigor:} Laboratory control with randomization
    \item \textbf{Replication:} Findings must be robust across contexts
    \item \textbf{Formal models:} Fehr-Schmidt inequality aversion has precise parameters $(\alpha, \beta)$
    \item \textbf{Falsifiable predictions:} Social preferences predict voting, crime, redistribution
    \item \textbf{Heterogeneity:} Discrete types, not infinitely flexible parameters
\end{enumerate}

\subsubsection{What Fehr Implicitly Requires}

From these commitments, we can infer:

\begin{tcolorbox}[colback=gray!10!white,colframe=gray!75!black]
\textbf{Implicit Fehr Principle:}

A model is scientifically valuable in proportion to what it \textit{excludes}, not what it \textit{includes}.
\end{tcolorbox}

The Fehr-Schmidt model is valuable precisely because:
\begin{itemize}[noitemsep]
    \item It has only two parameters ($\alpha$, $\beta$)
    \item These parameters have restricted ranges
    \item The model makes clear predictions that can fail
    \item Alternative behaviors would falsify the model
\end{itemize}

\subsection{The Problem with Integration}

\subsubsection{The Historical Achievement of Behavioral Economics}

The behavioral economics revolution integrated:
\begin{itemize}[noitemsep]
    \item Social preferences (Fehr \& Schmidt, 1999)
    \item Context-dependence (Kahneman \& Tversky)
    \item Norm-following (Bicchieri)
    \item Institutional effects (Ostrom)
    \item Cognitive limitations (bounded rationality)
\end{itemize}

\subsubsection{The Methodological Danger}

\begin{tcolorbox}[colback=red!5!white,colframe=red!75!black,title=\textbf{Fehr's Alarm}]
\begin{center}
\textbf{If everything can be integrated, nothing is excluded.}
\end{center}

When a framework can accommodate:
\begin{itemize}[noitemsep]
    \item Any preference structure
    \item Any context effect
    \item Any norm
    \item Any heterogeneity
\end{itemize}

Then every observation is consistent with the framework, and the framework is unfalsifiable.
\end{tcolorbox}

\subsection{Arrow-Debreu as Disciplinary Device}

\subsubsection{The Standard Interpretation}

Arrow-Debreu general equilibrium is typically criticized as:
\begin{itemize}[noitemsep]
    \item Unrealistic (perfect markets, complete information)
    \item Descriptively false (markets fail, preferences are social)
    \item Normatively problematic (assumes homo economicus)
\end{itemize}

\subsubsection{Fehr's Implicit Interpretation}

For Fehr, Arrow-Debreu serves a different function:

\begin{tcolorbox}[colback=blue!5!white,colframe=blue!75!black]
\textbf{Arrow-Debreu as Null Hypothesis:}

The value of Arrow-Debreu is not descriptive accuracy but \textit{methodological discipline}. It provides a precise null hypothesis against which deviations can be measured.
\end{tcolorbox}

This explains Fehr's research strategy:
\begin{enumerate}
    \item Start with the standard model (additive, selfish preferences)
    \item Design experiments that test specific deviations
    \item Measure the magnitude and conditions of deviation
    \item Propose minimal extensions (e.g., two parameters $\alpha$, $\beta$)
\end{enumerate}

\subsubsection{The Implicit Demand}

Fehr's implicit demand for complementarity frameworks:

\begin{quote}
``Show me first the case without complementarity. Then show me precisely which complementarity is necessary and sufficient. Do not give me a framework where complementarity is the starting assumption.''
\end{quote}


% ============================================================================
% PART II: WHY COMPLEMENTARITY IS MAXIMALLY THREATENING
% ============================================================================
\section{Why Complementarity Maximizes the Integration Problem}
\label{sec:complementarity-threat}

\subsection{The Structure of the Threat}

Complementarity ($\gamma$) is not merely one more element to integrate. It is a \textit{meta-principle} that explains how elements interact. This makes it uniquely threatening to falsifiability.

\subsubsection{Ordinary Integration}

Adding social preferences to standard economics:
\[
U = U^F + U^{SS}
\]
This is testable because:
\begin{itemize}[noitemsep]
    \item $U^{SS}$ has specific functional forms
    \item Parameters can be estimated
    \item The model fails if $U^{SS}$ has wrong structure
\end{itemize}

\subsubsection{Complementarity Integration}

Adding complementarity:
\[
U = U^F + U^{SS} + \gamma_{F,SS} \cdot U^F \cdot U^{SS}
\]

Now any deviation can be attributed to $\gamma$:
\begin{itemize}[noitemsep]
    \item If behavior exceeds prediction: ``positive complementarity''
    \item If behavior falls short: ``negative complementarity''
    \item If behavior varies by context: ``context-dependent $\gamma$''
\end{itemize}

\subsection{The Worst-Case Scenario: Universal Integration}

\begin{tcolorbox}[colback=red!10!white,colframe=red!75!black,title=\textbf{Fehr's Nightmare}]
\begin{center}
\textbf{A formal framework that can integrate any observation through complementarity parameters.}
\end{center}

In this scenario:
\begin{itemize}[noitemsep]
    \item Every model fits ex post
    \item No prediction can fail (just adjust $\gamma$)
    \item The framework explains everything, therefore nothing
\end{itemize}
\end{tcolorbox}

\subsection{The ``Missing Complementarity'' Escape}

The most dangerous pattern:

\begin{enumerate}
    \item Framework predicts outcome $Y^*$
    \item Actual outcome is $Y \neq Y^*$
    \item Analyst says: ``The complementarity $\gamma_{ab}$ was missing''
    \item Framework is ``saved'' by adding $\gamma_{ab}$
    \item This can repeat indefinitely
\end{enumerate}

This is precisely the pattern that Lakatos (1970) identified as \textit{degenerative} research programs: protecting the core by adding auxiliary hypotheses.

\subsection{Why This Is Not the Standard ``Too Many Parameters'' Critique}

The standard critique (Appendix CRT, Argument 6) concerns:
\begin{itemize}[noitemsep]
    \item Computational tractability
    \item Statistical power
    \item Estimation feasibility
\end{itemize}

Fehr's deeper concern is different:

\begin{tcolorbox}[colback=gray!10!white,colframe=gray!75!black]
\textbf{Fehr's Meta-Concern:}

Even if we could estimate all $\gamma$ parameters perfectly, a framework that accommodates any $\gamma$ pattern is not scientific.

The issue is not ``too many parameters'' but ``no excluded parameter patterns.''
\end{tcolorbox}


% ============================================================================
% PART III: THE RESOLUTION - GAMMA AS EXCLUSION PRINCIPLE
% ============================================================================
\section{The Resolution: $\gamma$ as Exclusion Principle}
\label{sec:resolution}

\subsection{Two Uses of Complementarity}

\begin{table}[h]
\centering
\caption{Integration vs. Exclusion Principle}
\label{tab:integration-exclusion}
\renewcommand{\arraystretch}{1.4}
\begin{tabular}{@{}p{6cm}p{6cm}@{}}
\toprule
\textbf{$\gamma$ as Integration Principle} & \textbf{$\gamma$ as Exclusion Principle} \\
\midrule
``Everything is complementary to some degree'' & ``Most things are \textit{not} complementary'' \\
$\gamma$ explains why things fit together & $\gamma = 0$ is the null hypothesis \\
Deviations = missing complementarity & Deviations = model failure \\
Post-hoc rationalization possible & Ex-ante predictions required \\
Unfalsifiable framework & Falsifiable claims \\
\bottomrule
\end{tabular}
\end{table}

\subsection{The Fehr-Compatible Formulation}

For EBF to satisfy Fehr's methodological requirements:

\begin{tcolorbox}[colback=green!5!white,colframe=green!75!black,title=\textbf{Fehr-Compatible EBF Principle}]
\begin{center}
\large
\textbf{The default is $\gamma = 0$. Every $\gamma \neq 0$ requires justification.}
\end{center}

\vspace{0.5em}
This means:
\begin{enumerate}
    \item Additivity is the null hypothesis
    \item Complementarity must be predicted ex ante
    \item Failed predictions falsify the $\gamma$ claim
    \item The framework tells us where $\gamma = 0$
\end{enumerate}
\end{tcolorbox}

\subsection{Formal Implementation}

\subsubsection{The Null Hypothesis Structure}

For any pair $(a, b)$:
\begin{equation}
H_0: \gamma_{ab} = 0 \quad \text{(additivity)}
\end{equation}
\begin{equation}
H_1: \gamma_{ab} \neq 0 \quad \text{(complementarity or substitutability)}
\end{equation}

The framework commits to:
\begin{itemize}[noitemsep]
    \item Specifying which $H_1$ claims it makes
    \item Specifying what evidence would falsify each claim
    \item Accepting falsification when evidence contradicts
\end{itemize}

\subsubsection{The Exclusion Zone}

For scientific legitimacy, EBF must specify:

\begin{definition}[Exclusion Zone]
The exclusion zone $\mathcal{E}$ is the set of $(a, b)$ pairs for which the framework commits to $\gamma_{ab} = 0$:
\[
\mathcal{E} = \{(a,b) : \gamma_{ab} = 0 \text{ by theoretical commitment}\}
\]
\end{definition}

The framework is falsified if $\gamma_{ab} \neq 0$ is found for $(a,b) \in \mathcal{E}$.

\subsection{Examples of Exclusion Commitments}

\begin{table}[h]
\centering
\caption{Example Exclusion Commitments}
\label{tab:exclusion-examples}
\renewcommand{\arraystretch}{1.3}
\begin{tabular}{@{}p{4cm}p{4cm}p{4cm}@{}}
\toprule
\textbf{Domain} & \textbf{Exclusion Claim} & \textbf{Falsification} \\
\midrule
Time preferences & $\gamma_{patience, risk} = 0$ for independent traits & Evidence of interaction \\
Social preferences & $\gamma_{\alpha, \beta}$ constrained by FS model & Unconstrained $(\alpha, \beta)$ \\
Context effects & Some $\Psi$ dimensions are orthogonal & Cross-effects observed \\
Institutional design & Minimal mechanism conditions & Violations of independence \\
\bottomrule
\end{tabular}
\end{table}

\subsection{The One Sentence for Fehr}

If Ernst Fehr required a single statement of EBF's position:

\begin{tcolorbox}[colback=blue!10!white,colframe=blue!75!black]
\begin{center}
\large
\textit{``Complementarity is for me not an integration principle, but an instrument for the exclusion of unnecessary mechanisms.''}
\end{center}
\end{tcolorbox}

Or, more sharply:

\begin{tcolorbox}[colback=blue!10!white,colframe=blue!75!black]
\begin{center}
\large
\textit{``My goal is not to show that everything is complementary, but precisely to show when nothing is complementary.''}
\end{center}
\end{tcolorbox}

These formulations are 100\% Fehr-compatible.


% ============================================================================
% PART IV: IMPLICATIONS FOR EBF METHODOLOGY
% ============================================================================
\section{Implications for EBF Methodology}
\label{sec:implications}

\subsection{Methodological Requirements}

To maintain scientific legitimacy under Fehr's critique, EBF must:

\begin{enumerate}
    \item \textbf{Ex-ante specification:} Every $\gamma$ claim must be stated before data analysis

    \item \textbf{Null hypothesis priority:} $\gamma = 0$ is always the default

    \item \textbf{Falsification criteria:} Every $\gamma$ claim specifies what would refute it

    \item \textbf{Exclusion documentation:} The framework explicitly states where $\gamma = 0$

    \item \textbf{No post-hoc rescue:} Failed predictions are documented, not explained away
\end{enumerate}

\subsection{Connection to Appendix SUN}

The ten predictions in Appendix SUN (PREDICT-FALSIFIABLE) exemplify this approach:

\begin{itemize}[noitemsep]
    \item Each prediction specifies direction and magnitude
    \item Each prediction states falsification criteria
    \item Predictions can fail, and failure is informative
\end{itemize}

Appendix SUN should be understood as the \textit{application} of the exclusion principle: it specifies what EBF predicts will \textit{not} happen (e.g., UCT effects will \textit{not} be equal across social capital levels).

\subsection{Connection to Appendix CRT}

The ten canonical arguments (Appendix CRT) concern \textit{technical} challenges:
\begin{itemize}[noitemsep]
    \item How to identify $\gamma$
    \item How to handle scale
    \item How to manage complexity
\end{itemize}

This appendix (XVIII) concerns the \textit{meta-level} challenge:
\begin{itemize}[noitemsep]
    \item Whether $\gamma$ frameworks can be scientific at all
    \item What makes them scientific (exclusion, not integration)
\end{itemize}

\subsection{The Research Program Test}

Following Lakatos (1970), a research program is \textit{progressive} if it:
\begin{enumerate}
    \item Predicts novel facts
    \item Some predictions are confirmed
    \item Failed predictions lead to learning, not ad-hoc rescue
\end{enumerate}

EBF passes this test when:
\begin{itemize}[noitemsep]
    \item $\gamma$ claims generate novel predictions (Appendix SUN)
    \item Predictions are tested and some confirmed
    \item Failed predictions are documented (Appendix EPS)
\end{itemize}

EBF fails this test when:
\begin{itemize}[noitemsep]
    \item $\gamma$ is invoked post-hoc to save failed predictions
    \item Every deviation is attributed to ``missing complementarity''
    \item The framework never excludes anything
\end{itemize}


% ============================================================================
% PART V: THE DEEPER PHILOSOPHICAL ISSUE
% ============================================================================
\section{The Deeper Philosophical Issue}
\label{sec:philosophy}

\subsection{Two Conceptions of Science}

The Fehr critique reveals a deeper tension between two conceptions of scientific progress:

\begin{table}[h]
\centering
\caption{Two Conceptions of Science}
\label{tab:two-conceptions}
\renewcommand{\arraystretch}{1.4}
\begin{tabular}{@{}p{6cm}p{6cm}@{}}
\toprule
\textbf{Unification View} & \textbf{Discipline View} \\
\midrule
Science seeks unified theories & Science seeks precise, testable models \\
Progress = more integration & Progress = more exclusion \\
A good theory explains everything & A good theory explains little, precisely \\
Paradigm: General Relativity & Paradigm: Molecular Biology \\
Hero: Einstein & Hero: Popper \\
\bottomrule
\end{tabular}
\end{table}

\subsection{Fehr's Position}

Fehr is clearly in the \textit{discipline} camp:
\begin{itemize}[noitemsep]
    \item He prefers models with few parameters
    \item He values experimental falsification
    \item He is suspicious of grand theoretical frameworks
    \item He emphasizes replication and robustness
\end{itemize}

\subsection{The EBF Response}

EBF can satisfy both views by recognizing:

\begin{tcolorbox}[colback=green!5!white,colframe=green!75!black]
\textbf{Synthesis:}

A unified framework is scientific if and only if it generates testable submodels that can fail independently.

EBF is the \textit{organizational structure} for testable models, not a substitute for them.
\end{tcolorbox}

This is why EBF is a \textit{framework}, not a \textit{model}:
\begin{itemize}[noitemsep]
    \item Frameworks organize knowledge
    \item Models make predictions
    \item Frameworks are judged by the models they generate
    \item Models are judged by their predictions
\end{itemize}

\subsection{The Framework-Model Distinction}

\begin{table}[h]
\centering
\caption{Framework vs. Model}
\label{tab:framework-model}
\renewcommand{\arraystretch}{1.3}
\begin{tabular}{@{}lll@{}}
\toprule
\textbf{Property} & \textbf{Framework (EBF)} & \textbf{Model (Fehr-Schmidt)} \\
\midrule
Parameters & Many, organized & Few, estimated \\
Predictions & Via derived models & Direct \\
Falsifiability & Via generated models & Direct \\
Scope & Universal & Specific \\
Scientific status & Organizing structure & Testable hypothesis \\
\bottomrule
\end{tabular}
\end{table}

EBF is to specific models as:
\begin{itemize}[noitemsep]
    \item The periodic table is to chemical reactions
    \item Evolutionary theory is to specific phylogenies
    \item Game theory is to specific games
\end{itemize}


% ============================================================================
% PART VI: TAXONOMY OF INTEGRABILITY
% ============================================================================
\section{Taxonomy: Integrable vs. Non-Integrable Models}
\label{sec:taxonomy}

This taxonomy is not ontological but methodological: it does not say what is true,
but what may be modeled together without loss of falsifiability.

\subsection{Integrable Models (Under Clear Conditions)}

These models can be combined when the integration conditions are explicit.

\subsubsection{Class A: Additive Baseline + Single Deviation}

\begin{itemize}[noitemsep]
    \item Arrow-Debreu + inequality aversion
    \item Rational choice + bounded rationality (local)
\end{itemize}

\textbf{Integrable because:}
\begin{itemize}[noitemsep]
    \item Clear null hypothesis
    \item Clear deviation
    \item Clear testability
\end{itemize}

\subsubsection{Class B: Complementarity Within a Single Mechanism}

\begin{itemize}[noitemsep]
    \item Incentive $\times$ Information
    \item Norm $\times$ Expectation
\end{itemize}

\textbf{Integrable if:}
\begin{itemize}[noitemsep]
    \item Mechanism is identical
    \item Interaction is specified ex ante
\end{itemize}

\subsubsection{Class C: Temporally Sequenced Models}

\begin{itemize}[noitemsep]
    \item Short-term: Incentives
    \item Long-term: Norm dynamics
\end{itemize}

\textbf{Integrable if:}
\begin{itemize}[noitemsep]
    \item Temporal structure is explicit
    \item No simultaneity is claimed
\end{itemize}

\subsection{Non-Integrable Models (Exclusion Required)}

This is Fehr's critical point: some models \textit{must} be excluded.

\begin{table}[h]
\centering
\caption{Non-Integrable Model Pairs}
\label{tab:non-integrable}
\renewcommand{\arraystretch}{1.4}
\small
\begin{tabular}{@{}p{4cm}p{4cm}p{4cm}@{}}
\toprule
\textbf{Model A} & \textbf{Model B} & \textbf{Why Non-Integrable} \\
\midrule
Rational expectations & Systematic misperception & Errors are either random or systematic \\
Stable preferences & Constructed preferences & ``Preference'' becomes meaningless \\
Inequality aversion & Reciprocity & Identification collapses on same data \\
Pareto welfare & Rawlsian welfare & Normative choice required, not technique \\
Institutions explain behavior & Behavior explains institutions & Circularity without temporal structure \\
\bottomrule
\end{tabular}
\end{table}

\begin{tcolorbox}[colback=red!5!white,colframe=red!75!black,title=\textbf{Fehr-Compatible Principle}]
\begin{center}
\textbf{A model is good precisely where it excludes other models.}
\end{center}
\end{tcolorbox}


% ============================================================================
% PART VII: WORKED EXAMPLE
% ============================================================================
\section{Worked Example: Non-Integrative Paper Design}
\label{sec:worked-example}

This section provides a concrete blueprint for a paper that Fehr would accept.

\subsection{Example Title}

\begin{quote}
\textit{``When Complementarity Fails: Evidence from a No-Interaction Benchmark''}
\end{quote}

\subsection{Structure}

\subsubsection{Step 1: Disciplinary Starting Point}

\begin{itemize}[noitemsep]
    \item Arrow-Debreu baseline
    \item Additive preferences
    \item No complementarity
\end{itemize}

\textbf{Explicit null hypothesis:} $\gamma = 0$

\subsubsection{Step 2: Single Deviation}

Introduce exactly one interaction, e.g., Incentive $\times$ Norm.

\textbf{Explicitly excluded (not identifiable in design):}
\begin{itemize}[noitemsep]
    \item Social image effects
    \item Reciprocity
    \item Learning models
\end{itemize}

\subsubsection{Step 3: Experimental Design}

\begin{itemize}[noitemsep]
    \item Factorial structure
    \item Clear counterfactuals
    \item Measurement: main effects + interaction
\end{itemize}

\subsubsection{Step 4: Central Tests}

\begin{align}
H_0&: \gamma = 0 \quad \text{(additivity)} \\
H_1&: \gamma \neq 0 \quad \text{(complementarity)}
\end{align}

\textbf{Critical:} If $\gamma = 0$, the model \textit{fails}---no rescue through ``missing integration.''

\subsubsection{Step 5: Discussion (Decisive)}

\begin{itemize}[noitemsep]
    \item Where does the model \textit{not} explain?
    \item Which alternative models remain \textit{excluded}?
    \item Why are they \textit{not integrable}?
\end{itemize}

This is where the paper gains scientific value.

\subsubsection{Step 6: Conclusion}

\begin{quote}
\textit{Complementarity is not a default assumption, but a risky hypothesis.}
\end{quote}

This is 100\% Fehr-compatible.


% ============================================================================
% PART VIII: METHODOLOGICAL BOX TEMPLATE
% ============================================================================
\section{Template: Methodological Box for Manuscripts}
\label{sec:template-box}

The following box can be used verbatim in EBF-derived papers.

\begin{tcolorbox}[colback=gray!10!white,colframe=gray!75!black,
    title=\textbf{Box: Why We Do Not Integrate All Models}]

In this paper, we deliberately refrain from integrating multiple behavioral mechanisms.
While modern frameworks allow the formal combination of incentives, norms, preferences,
beliefs, and institutional context, such integration often eliminates falsifiability.

Our approach treats complementarity not as an integrative principle, but as a selective one.
We begin with a restrictive benchmark in which no complementarity is present and introduce
interaction effects only where they generate testable deviations.

Models that are logically, normatively, or identificationally incompatible with this
structure are explicitly excluded. This exclusion is not a limitation but a methodological
necessity: \textbf{a model that can accommodate all mechanisms explains none.}

Complementarity, in our framework, is therefore a hypothesis to be rejected,
not a narrative to be confirmed.

\end{tcolorbox}

\subsection{When to Use This Box}

Include this box (or a variant) when:
\begin{enumerate}
    \item Your paper uses EBF or $\gamma$-framework concepts
    \item Reviewers might expect integration of multiple mechanisms
    \item You need to justify why certain models are excluded
    \item You want to signal methodological rigor to Fehr-tradition reviewers
\end{enumerate}

\subsection{The Summary Principle}

\begin{tcolorbox}[colback=blue!5!white,colframe=blue!75!black]
\begin{center}
\large
\textbf{Complementarity is scientifically valuable only when it forces us to say which models do not apply.}
\end{center}
\end{tcolorbox}

This is where you take Fehr seriously---and go beyond him.


% ============================================================================
% PART IX: EXPLAINABILITY VS. COMBINABILITY
% ============================================================================
\section{The Core Distinction: Explainability vs. Combinability}
\label{sec:explainability-combinability}

This section formalizes the central methodological distinction that underlies
the entire Fehr critique.

\subsection{The Precise Thesis}

\begin{tcolorbox}[colback=red!5!white,colframe=red!75!black,title=\textbf{Central Claim}]
\begin{center}
\large
\textbf{EBF can explain any single model---but it cannot combine these models without losing its scientific discipline.}
\end{center}

\vspace{0.5em}
This is not a contradiction but a structural property of evidence-based approaches.
\end{tcolorbox}

\subsection{Why EBF Can Explain Any Model}

Evidence-based behavioral approaches operate as follows:
\begin{itemize}[noitemsep]
    \item They ask: ``Does intervention $X$ work under condition $Y$?''
    \item They require: a causal effect in a clearly defined context
    \item They do \textit{not} commit to a comprehensive theory
\end{itemize}

Therefore, EBF can legitimately say:
\begin{itemize}[noitemsep]
    \item This result is consistent with inequality aversion
    \item This result is consistent with reciprocity
    \item This result is consistent with social norms
    \item This result is consistent with incentive effects
\end{itemize}

EBF is:
\begin{itemize}[noitemsep]
    \item \textbf{Local:} context-specific findings
    \item \textbf{Post hoc:} explains observed effects
    \item \textbf{Instrumental:} focuses on what works
    \item \textbf{Agnostic:} does not adjudicate which model is ``true''
\end{itemize}

\subsection{Why EBF Cannot Combine These Models}

\textbf{Combination} means:
\begin{itemize}[noitemsep]
    \item Multiple mechanisms simultaneously
    \item In a unified model
    \item With shared parameters
    \item With shared counterfactuals
\end{itemize}

This fails systematically for four reasons:

\subsubsection{(a) Identification Collapse}

The same data can be explained by multiple mechanisms.
No clean separation is possible.

\[
\text{Combination} \Rightarrow \text{Non-Falsifiability}
\]

\subsubsection{(b) Logical Contradictions}

Example: \textit{stable preferences} and \textit{constructed preferences} cannot both be true.
Combining them makes ``preference'' meaningless.

\subsubsection{(c) Hidden Normative Decisions}

\begin{itemize}[noitemsep]
    \item Which mechanism counts when?
    \item Which welfare function applies?
\end{itemize}

EBF cannot resolve this neutrally---it hides the normative choice.

\subsubsection{(d) Post-Hoc Rationalization}

When combination is allowed, every failure can be explained:

\begin{quote}
``It didn't work because we didn't integrate component $Z$.''
\end{quote}

This is precisely Fehr's nightmare.

\subsection{The Central Methodological Distinction}

\begin{table}[h]
\centering
\caption{Explainability vs. Combinability}
\label{tab:explain-combine}
\renewcommand{\arraystretch}{1.4}
\begin{tabular}{@{}p{5.5cm}p{5.5cm}@{}}
\toprule
\textbf{Explainability (EBF)} & \textbf{Combinability (Theory)} \\
\midrule
Local & Global \\
Post hoc & Ex ante \\
Context-bound & Restriktive \\
Pluralistic (many models) & Falsifiable (one model) \\
Asks: ``What works?'' & Asks: ``What's true?'' \\
\bottomrule
\end{tabular}
\end{table}

\textbf{EBF deliberately operates only at the explainability level.}

\subsection{Arrow-Debreu Makes the Difference Visible}

Arrow-Debreu general equilibrium:
\begin{itemize}[noitemsep]
    \item Is not realistic
    \item But is highly restrictive
    \item Shows what happens when no complementarity exists
\end{itemize}

EBF, by contrast:
\begin{itemize}[noitemsep]
    \item Departs from the AD framework
    \item Without setting a new disciplining benchmark
\end{itemize}

Therefore:
\begin{itemize}[noitemsep]
    \item EBF can explain everything
    \item But EBF can exclude nothing
\end{itemize}

\subsubsection{Critical Clarification: EBF Does Not Explain Arrow-Debreu}

This point requires precision. Fehr's implicit claim is:

\begin{tcolorbox}[colback=yellow!5!white,colframe=yellow!75!black,title=\textbf{What EBF Actually Does (and Does Not Do)}]
\textbf{EBF does not explain the Arrow-Debreu model.}

EBF merely identifies contexts in which specific interaction terms are empirically negligible.
\end{tcolorbox}

Why this distinction matters:

\textbf{(a) AD sets \textit{all} relevant complementarities to zero simultaneously.}\\
EBF sets \textit{individual} interactions to zero \textit{locally}.\\
These are logically not the same.

\textbf{(b) AD is globally consistent} (all markets, all agents, all decisions).\\
EBF works with partial equilibria, local interventions, limited horizons.

\textbf{(c) AD is ex ante disciplining.}\\
EBF is ex post explanatory.

AD says: ``If this structure holds, then equilibrium $X$ follows.''\\
EBF says: ``Here we see effect $Y$, so this interaction was probably not relevant.''

This is not an explanation of AD, but a \textit{post-hoc classification of contexts}.

\subsubsection{The Precise Logical Relationship}

\begin{quote}
\textit{EBF can show that certain real situations behave \textbf{as if} some complementarities were zero.}

\textit{But EBF cannot show that the world corresponds to the Arrow-Debreu model.}
\end{quote}

Because:
\begin{itemize}[noitemsep]
    \item AD is a \textbf{construction principle}
    \item EBF is a \textbf{measurement procedure}
\end{itemize}

\subsubsection{Why Fehr Insists on This Distinction}

Without it, the following error occurs:
\begin{enumerate}
    \item Observe an isolated effect
    \item Say: ``There was apparently no complementarity here''
    \item Act as if AD has been explained or confirmed
\end{enumerate}

Fehr's response:

\begin{quote}
\textit{No. You have only shown that in this case, interaction played no role. That is much weaker---and must be named as such.}
\end{quote}

\subsubsection{The Deeper Methodological Point}

Fehr protects something central here:

\begin{tcolorbox}[colback=gray!10!white,colframe=gray!75!black]
\textbf{The right of theories to be more than an aggregation of empirical findings.}
\end{tcolorbox}

AD is:
\begin{itemize}[noitemsep]
    \item Not ``extrapolated'' from data
    \item Not reconstructible through EBF
    \item A deliberately unrealistic ordering framework
\end{itemize}

EBF cannot replace this framework---it can only work within or alongside it.

\subsubsection{The Clean Formula}

\begin{tcolorbox}[colback=green!5!white,colframe=green!75!black]
\begin{center}
\textit{``Evidence-based approaches do not explain the Arrow-Debreu model.\\
They merely identify contexts in which specific interaction terms are empirically negligible.''}
\end{center}
\end{tcolorbox}

Or equivalently:

\begin{quote}
\textit{Arrow-Debreu is a theoretical reference point.\\
EBF is an empirical classification instrument.\\
Complementarity connects both meaningfully only when it is explicitly localized and bounded.}
\end{quote}

\subsection{Where Fehr Is Entirely Correct}

Fehr's implicit position:

\begin{tcolorbox}[colback=gray!10!white,colframe=gray!75!black]
\textit{As long as EBF does not say which models do not apply, it is not a theory but a catalogue of working interventions.}
\end{tcolorbox}

This is not a criticism but a boundary demarcation.

\subsection{The Correct Conclusion}

The right statement is \textbf{not}:

\begin{quote}
``EBF is bad.''
\end{quote}

But rather:

\begin{quote}
\textbf{``EBF is deliberately non-combinatorial.''}
\end{quote}

EBF is:
\begin{itemize}[noitemsep]
    \item Evidence-strong
    \item Policy-proximate
    \item Practical
\end{itemize}

But EBF is \textbf{not}:
\begin{itemize}[noitemsep]
    \item Theoretically integrating
    \item Paradigmatically disciplining
\end{itemize}

\subsection{The Gap That Complementarity Theory Fills}

Your work addresses exactly this gap:

\begin{tcolorbox}[colback=blue!5!white,colframe=blue!75!black]
\textbf{Not} explaining which model fits,\\
\textbf{but} explaining when models \textit{cannot} jointly hold.
\end{tcolorbox}

Or in one sentence:

\begin{quote}
\large
\textit{EBF explains effects; Complementarity decides between models.}
\end{quote}

\subsection{The Exact Formulation for Academic Use}

The following sentence is correct, sharp, Fehr-compatible, and non-polemical:

\begin{tcolorbox}[colback=green!5!white,colframe=green!75!black]
\begin{center}
\large
\textit{``Evidence-based approaches can rationalize virtually any behavioral model in isolation, but they lose their scientific discipline once these models are combined.''}
\end{center}
\end{tcolorbox}

This is where a disciplined complementarity theory begins.


% ============================================================================
% SUMMARY
% ============================================================================
\section{Summary}
\label{sec:summary}

\begin{tcolorbox}[colback=blue!5!white,colframe=blue!75!black,title=\textbf{Concluding Statement}]

\textbf{Fehr's Real Problem:}

Not that complementarity is false, but that complementarity frameworks threaten the disciplinary power that makes models scientific.

\textbf{The Resolution:}

EBF is scientifically legitimate when $\gamma$ functions as an \textit{exclusion principle}:
\begin{itemize}[noitemsep]
    \item $\gamma = 0$ is the null hypothesis
    \item $\gamma \neq 0$ requires ex-ante justification
    \item Failed $\gamma$ predictions falsify claims
    \item The framework specifies where complementarity does \textit{not} hold
\end{itemize}

\textbf{The One Sentence:}

\textit{``Complementarity is an instrument for exclusion, not integration.''}

\end{tcolorbox}

\subsection{What This Means---Concretely}

\subsubsection{For EBF}

EBF is methodologically correct but deliberately limited.

\textbf{EBF says:}
\begin{itemize}[noitemsep]
    \item ``This intervention had an effect here.''
    \item ``This behavior is consistent with Model $X$.''
\end{itemize}

\textbf{EBF does not say:}
\begin{itemize}[noitemsep]
    \item ``This model is universally valid.''
    \item ``These models can be combined.''
\end{itemize}

\textbf{Consequence:} EBF is an evidence filter, not a theory construction.
If EBF is used as theory anyway, Fehr's nightmare occurs.

\subsubsection{For Complementarity}

Complementarity must \textbf{not} mean: ``Everything is connected to everything.''

Complementarity \textbf{must} mean: ``The additive benchmark breaks here---and only here.''

Concretely:
\begin{itemize}[noitemsep]
    \item Complementarity is not a default assumption
    \item It is a risky deviation hypothesis
    \item It must be able to fail
\end{itemize}

\textbf{Consequence:} Complementarity is legitimate only as exception, not as world explanation.

\subsubsection{For Models}

\textbf{Correct:}
\begin{itemize}[noitemsep]
    \item One model per explanation
    \item Clear null hypothesis
    \item Clear exclusions
\end{itemize}

\textbf{Incorrect:}
\begin{itemize}[noitemsep]
    \item Hybrid models containing everything
    \item ``Integration frameworks''
    \item Meta-models without exclusion rules
\end{itemize}

\textbf{Consequence:} Models must be able to exclude each other.

\subsubsection{For Arrow-Debreu}

Arrow-Debreu is not refuted but methodologically irreplaceable.

AD shows: \textit{what happens when there is no complementarity.}

This is:
\begin{itemize}[noitemsep]
    \item Not realism
    \item But a disciplining anchor
\end{itemize}

\textbf{Consequence:} Without AD (or something similar), complementarity becomes arbitrary.

\subsubsection{For Fehr}

Fehr's position is not anti-behavioral, but:

\begin{quote}
\textit{``If you can integrate everything, you must first say what you do \textbf{not} integrate.''}
\end{quote}

He accepts EBF, Behavioral Economics, even complementarity.\\
He rejects: ``everything-explains-everything'' models.

\textbf{Consequence:} His point is methodological honesty.

\subsection{The Operative Consequence (Practical!)}

When writing a paper, experiment, or policy analysis today:

\begin{enumerate}
    \item \textbf{Start with a null model}\\
    (e.g., Arrow-Debreu / additive preferences)

    \item \textbf{Introduce one deviation}\\
    (e.g., a specific complementarity)

    \item \textbf{Explicitly exclude other models}\\
    (not from laziness, but from logic)

    \item \textbf{Accept failure}\\
    (if $\gamma = 0$, that is the result)
\end{enumerate}

\textbf{This is science.}

\subsection{The Final Formulations}

\begin{tcolorbox}[colback=green!5!white,colframe=green!75!black,title=\textbf{Memorize These}]

\textbf{Statement 1:}\\
\textit{EBF explains effects.\\
Theory decides between models.\\
Complementarity is only theory when it can also say ``No.''}

\vspace{1em}
\textbf{Statement 2 (simplest form):}\\
\textit{We may observe everything---\\
but we may not believe everything simultaneously.}

\end{tcolorbox}

\subsection{Why This Distinction Is Existential for EBF}

The statement ``EBF does not explain Arrow-Debreu, but merely identifies contexts where
specific complementarities are empirically negligible'' sounds technical. But it is the
\textbf{methodological protection mechanism} that keeps EBF scientifically legitimate.

\subsubsection{What Goes Wrong If We Ignore This Distinction}

\textbf{(1) EBF becomes ``Anything Goes''}

If we say ``EBF shows that AD is false'' or ``EBF replaces AD,'' then implicitly:
\begin{itemize}[noitemsep]
    \item Every observed behavior justifies a different model
    \item Without a common reference point
\end{itemize}
$\Rightarrow$ EBF explains everything---and therefore nothing.

\textbf{(2) Complementarity becomes an excuse}

Without AD as reference:
\begin{itemize}[noitemsep]
    \item Every intervention failure is explained by: ``We were missing another complementary measure''
\end{itemize}
$\Rightarrow$ Complementarity becomes unfalsifiable---and therefore unscientific.

\textbf{(3) Normative decisions become hidden}

If EBF ``replaces theory'':
\begin{itemize}[noitemsep]
    \item Policy decisions appear technically necessary, evidence-based, without alternatives
    \item But which interactions to consider \textit{is} a normative choice
\end{itemize}
$\Rightarrow$ EBF becomes technocratic.

\textbf{(4) Learning capacity is lost}

Without a null model:
\begin{itemize}[noitemsep]
    \item We don't know when something doesn't work
    \item Every result is explainable
\end{itemize}
$\Rightarrow$ Failure becomes uninformative. This is fatal for science.

\subsubsection{What EBF Can Only Be If We Follow This Distinction}

\textbf{EBF as classification instrument:}
\begin{quote}
``In this context, $\gamma_{AB} \approx 0$.''\\
``In that context, $\gamma_{AB}$ is significant.''\\
``For other interactions, we have no statement.''
\end{quote}
This is real evidence.

\textbf{EBF as bridge, not replacement:}
\begin{itemize}[noitemsep]
    \item EBF connects theory with practice---without replacing theory
    \item Theory defines null models and limits interpretations
\end{itemize}
Both remain necessary.

\textbf{Complementarity as hypothesis, not worldview:}
\begin{itemize}[noitemsep]
    \item Tested, rejected, or accepted
    \item Not presupposed, totalized, or ideologized
\end{itemize}
This is how it stays scientific.

\subsubsection{The Hardest But Most Honest Statement}

\begin{tcolorbox}[colback=red!10!white,colframe=red!75!black,title=\textbf{Final Warning}]
\begin{center}
\large
\textit{Without this distinction, EBF becomes unfalsifiable---\\
a method to justify any policy measure post hoc.}
\end{center}

\vspace{0.5em}
This is:
\begin{itemize}[noitemsep]
    \item Scientifically unacceptable
    \item Politically dangerous
    \item Precisely what Fehr warns against
\end{itemize}
\end{tcolorbox}

\subsubsection{The Complete Logic in One Sentence}

\begin{tcolorbox}[colback=blue!5!white,colframe=blue!75!black]
\textit{EBF is scientifically legitimate only when it does not claim to replace theories, but clearly states where empirically specific interactions play no role.}
\end{tcolorbox}


% ============================================================================
% THE PHILOSOPHY OF SCIENCE POSITION
% ============================================================================
\subsection{The Philosophy of Science Position}
\label{sec:philosophy-position}

This is the final level of abstraction. When we speak about the relationship between EBF, theories, and evidence, we must be clear about what \textit{kind} of knowledge each represents.

\subsubsection{Three Levels of Scientific Knowledge}

Scientific knowledge operates at three fundamentally distinct levels:

\begin{tcolorbox}[colback=blue!10!white,colframe=blue!75!black,title=\textbf{Level 1: Theory (Trans-Empirical)}]
\textbf{Function:} Structures the space of possibilities\\
\textbf{Characteristic:} Not derivable from data; defines what \textit{could} happen\\
\textbf{Example:} Arrow-Debreu general equilibrium, quantum mechanics, evolutionary theory\\
\textbf{Status:} Logically prior to observation; framework for interpretation
\end{tcolorbox}

\begin{tcolorbox}[colback=green!10!white,colframe=green!75!black,title=\textbf{Level 2: Empirical Evidence (Context-Bound)}]
\textbf{Function:} Locates deviations from theoretical predictions\\
\textbf{Characteristic:} Always contextual; reveals where theory does not apply\\
\textbf{Example:} Loss aversion in domain X, social preferences in population Y\\
\textbf{Status:} Cannot replace theory; can only constrain its application domain
\end{tcolorbox}

\begin{tcolorbox}[colback=yellow!10!white,colframe=yellow!75!black,title=\textbf{Level 3: Model Choice (Methodological-Normative)}]
\textbf{Function:} Decides which theory to apply in which context\\
\textbf{Characteristic:} Not purely empirical; involves normative judgment\\
\textbf{Example:} ``Use AD for policy X, behavioral model for intervention Y''\\
\textbf{Status:} Requires both theory and evidence, but is reducible to neither
\end{tcolorbox}

\subsubsection{The Category Error}

The deepest error in discussions of EBF is \textbf{category confusion}---treating empirical findings as if they could replace theoretical structures.

\textbf{The error manifests as:}
\begin{itemize}[noitemsep]
    \item ``EBF shows that AD is wrong'' (reduces theory to empirics)
    \item ``Behavioral economics replaces classical economics'' (conflates levels)
    \item ``Evidence determines which model is true'' (ignores Level 3)
\end{itemize}

\textbf{Why this is an error:}
\begin{enumerate}[noitemsep]
    \item Theories are not hypotheses that evidence can falsify directly
    \item Theories structure the \textit{interpretation} of evidence
    \item Evidence reveals where theories cease to apply, not that they are ``false''
    \item Model choice is a \textit{decision}, not a \textit{discovery}
\end{enumerate}

\subsubsection{The Historical Parallel: Why Positivism Failed}

This category error has a historical precedent: \textbf{logical positivism}.

\textbf{The positivist thesis:}
\begin{quote}
``All meaningful statements are either analytic (true by definition) or empirically verifiable.''
\end{quote}

\textbf{Why it failed:}
\begin{itemize}[noitemsep]
    \item The thesis itself is neither analytic nor empirically verifiable
    \item Theoretical terms (``electron'', ``utility'') cannot be reduced to observations
    \item Scientific practice requires trans-empirical commitments
\end{itemize}

\textbf{The parallel to EBF:}
\begin{itemize}[noitemsep]
    \item Claiming ``evidence shows AD is wrong'' repeats the positivist error
    \item Theoretical frameworks cannot be replaced by empirical patterns
    \item EBF, like any empirical program, \textit{presupposes} theoretical structure
\end{itemize}

Similarly, \textbf{behaviorism} failed for the same reason:
\begin{quote}
``Psychology should only study observable behavior, not mental states.''
\end{quote}
This proved impossible---behavior is interpretable only via theoretical constructs about mental states.

\subsubsection{The Role of Null Models}

Theories provide what evidence cannot: \textbf{null models}.

\begin{tcolorbox}[colback=gray!10!white,colframe=gray!75!black,title=\textbf{Why Null Models Are Trans-Empirical}]
A null model defines what would happen \textit{if} certain effects were absent.

\textbf{Example:} Arrow-Debreu defines equilibrium \textit{if} all agents are rational, all information is symmetric, all markets are complete.

This is not an empirical claim. It is a \textbf{construction principle} that enables:
\begin{itemize}[noitemsep]
    \item Identification of deviations
    \item Measurement of effect sizes
    \item Counterfactual reasoning
\end{itemize}

Without it, we cannot even \textit{define} what counts as a deviation.
\end{tcolorbox}

\textbf{EBF's proper relationship to null models:}
\begin{itemize}[noitemsep]
    \item EBF \textit{uses} null models (from theory)
    \item EBF \textit{measures} deviations from null models
    \item EBF \textit{does not provide} null models
    \item EBF \textit{cannot replace} null models
\end{itemize}

\subsubsection{The Canonical Sentence}

All of this condenses into a single formulation:

\begin{tcolorbox}[colback=red!5!white,colframe=red!75!black,title=\textbf{The Philosophy of Science Position}]
\begin{center}
\Large
\textbf{English:}\\[0.3em]
\textit{``Evidence does not explain theory;\\
it only reveals where theory ceases to apply.''}\\[1em]

\textbf{German:}\\[0.3em]
\textit{``Evidenz erklärt Theorie nicht;\\
sie zeigt nur, wo Theorie aufhört zu gelten.''}
\end{center}
\end{tcolorbox}

\textbf{Unpacking this sentence:}
\begin{enumerate}[noitemsep]
    \item \textbf{``Evidence does not explain theory''}: Empirical findings cannot replace theoretical structures
    \item \textbf{``it only reveals''}: Evidence has a locating function, not a replacing function
    \item \textbf{``where theory ceases to apply''}: Evidence defines application boundaries, not truth values
\end{enumerate}

\subsubsection{Implications for EBF Practice}

From the philosophy of science position, concrete implications follow:

\begin{table}[h!]
\centering
\begin{tabular}{p{0.45\textwidth}|p{0.45\textwidth}}
\toprule
\textbf{Wrong (Category Confusion)} & \textbf{Right (Level Distinction)} \\
\midrule
``EBF shows AD is false'' & ``EBF identifies contexts where AD predictions deviate'' \\
\midrule
``Behavioral economics replaces classical economics'' & ``Behavioral economics maps application boundaries of classical models'' \\
\midrule
``Evidence determines which model is true'' & ``Evidence informs model choice; choice remains normative'' \\
\midrule
``Complementarity explains why effects interact'' & ``Complementarity measures interaction strength; theory explains mechanism'' \\
\midrule
``EBF is a new theory of behavior'' & ``EBF is a measurement procedure that presupposes theoretical structure'' \\
\bottomrule
\end{tabular}
\caption{Category confusion vs. proper level distinction}
\end{table}

\subsubsection{The Three Foundational Traditions}

This position is grounded in three major philosophy of science traditions, each contributing a distinct insight.

\paragraph{Tradition I: Popper---Falsification and Null Models}

Popper's fundamental principle:
\begin{quote}
\textit{Theories are not derivable from data; they are conjectures that fail against data.}
\end{quote}

\textbf{Theories} (in the Popperian sense):
\begin{itemize}[noitemsep]
    \item Are conjectural
    \item Structure expectations
    \item Define what would constitute a counterexample
\end{itemize}

\textbf{Evidence} (in the Popperian sense):
\begin{itemize}[noitemsep]
    \item Tests theories
    \item Refutes (but does not replace) theories
    \item Cannot derive theories inductively
\end{itemize}

\textbf{Arrow-Debreu in Popperian terms:}

AD is \textit{not} an empirical claim. It is a \textbf{counterfactual null model}:
\begin{quote}
\textit{If there were no complementarities, no frictions, no interactions, then...}
\end{quote}

Such a statement cannot logically be ``explained.'' It can only:
\begin{itemize}[noitemsep]
    \item Serve as a reference point
    \item Be bounded by observed deviations
\end{itemize}

\textbf{EBF in Popperian terms:}

EBF is strictly falsificationist, local, and hypothetical. It says:
\begin{itemize}[noitemsep]
    \item \textit{Here} this assumption fails
    \item \textit{Here} this interaction is empirically $\approx 0$
\end{itemize}

This is fully Popperian---\textit{but only if} EBF does not claim to replace theory.

\begin{tcolorbox}[colback=red!5!white,colframe=red!75!black,title=\textbf{Popper's Warning (Directly Relevant)}]
Saying ``EBF explains AD'' commits exactly the error Popper warned against:
\begin{itemize}[noitemsep]
    \item Evidence is inductively elevated
    \item Theory is reduced to description
\end{itemize}
$\Rightarrow$ \textbf{Unscientific.}
\end{tcolorbox}

\paragraph{Tradition II: Lakatos---Research Programs and Protective Belts}

Lakatos extends Popper with a more nuanced structure.

\textbf{Research programs} (Lakatos) have:
\begin{itemize}[noitemsep]
    \item A \textbf{hard core}: never abandoned directly
    \item A \textbf{protective belt}: adjustable auxiliary hypotheses
\end{itemize}

\textbf{Arrow-Debreu = Hard Core}

AD is the hard core not because it is ``true,'' but because it:
\begin{itemize}[noitemsep]
    \item Disciplines
    \item Structures
    \item Enables comparison
\end{itemize}

The hard core states: \textit{Without interaction, without complementarity, without frictions...}

\textbf{EBF = Protective Belt}

EBF operates in the protective belt:
\begin{itemize}[noitemsep]
    \item Modifies assumptions
    \item Adds local deviations
    \item \textit{Without} replacing the core
\end{itemize}

EBF says: \textit{Here} this assumption is violated; \textit{here} we can set it to 0.

This is a \textbf{progressive protective belt}---exactly what good science does.

\textbf{What Fehr prevents (in Lakatosian terms):}

A \textbf{degenerative program}:
\begin{itemize}[noitemsep]
    \item Every deviation is explained
    \item Nothing is excluded
    \item Everything is integrated
\end{itemize}

This is exactly Fehr's critique of integrative behavioral economics and EBF misuse.

\begin{tcolorbox}[colback=orange!5!white,colframe=orange!75!black,title=\textbf{Fehr's Position in Lakatos's Language}]
\textit{EBF may extend the protective belt, but it must not dissolve the hard core.}

Otherwise:
\begin{itemize}[noitemsep]
    \item No discipline
    \item No progression
    \item Only accommodation
\end{itemize}
\end{tcolorbox}

\paragraph{Tradition III: Nancy Cartwright---Local Causality and the Dappled World}

Cartwright is decisive because she \textit{saves} EBF without sacrificing theory.

\textbf{Cartwright's core insight:}
\begin{quote}
\textit{The world is dappled. Laws do not hold everywhere.}
\end{quote}

There are no universal causal laws, only \textbf{local, context-bound efficacies}.

\textbf{EBF in Cartwright's terms:}

EBF is perfectly suited to:
\begin{itemize}[noitemsep]
    \item Identify local causalities
    \item State ``what works here''
\end{itemize}

EBF is \textit{not} suited to:
\begin{itemize}[noitemsep]
    \item Ground global theories
    \item Derive universal structures
\end{itemize}

This is why EBF is \textit{honest}---if it limits itself.

\textbf{Arrow-Debreu in Cartwright's terms:}

AD is not a ``law.'' It is a \textbf{theoretical scaffold}---an ordering framework.

EBF can show where this scaffold does not grip. EBF \textit{cannot} show:
\begin{itemize}[noitemsep]
    \item Why the scaffold exists
    \item That it has been ``explained''
\end{itemize}

\begin{tcolorbox}[colback=yellow!5!white,colframe=yellow!75!black,title=\textbf{Cartwright's Implicit Warning}]
If EBF overreaches:
\begin{itemize}[noitemsep]
    \item Local truths are globalized
    \item Policy becomes technocratic
\end{itemize}
$\Rightarrow$ Exactly what Fehr fears.
\end{tcolorbox}

\subsubsection{The Synthesis Across Three Traditions}

Each tradition contributes a distinct protective function:

\begin{table}[h!]
\centering
\begin{tabular}{p{0.2\textwidth}|p{0.35\textwidth}|p{0.35\textwidth}}
\toprule
\textbf{Tradition} & \textbf{AD's Role} & \textbf{EBF's Role} \\
\midrule
Popper & Counterfactual null model & Falsificationist testing \\
\midrule
Lakatos & Hard core of research program & Progressive protective belt \\
\midrule
Cartwright & Theoretical scaffold & Local causal identification \\
\bottomrule
\end{tabular}
\caption{The three traditions and their implications for AD and EBF}
\end{table}

\begin{tcolorbox}[colback=blue!5!white,colframe=blue!75!black,title=\textbf{The Canonical Formulation (Integrating All Three)}]
\begin{center}
\large
\textit{EBF does not explain theories (Popper),\\
it modifies protective belts (Lakatos),\\
and identifies local causality (Cartwright).}\\[0.5em]
\normalsize
\textbf{Anything else is a category error.}
\end{center}
\end{tcolorbox}

\subsubsection{What Happens If We Violate These Principles}

From each tradition, a specific failure mode follows:

\begin{enumerate}
    \item \textbf{Popper violation:} Loss of falsifiability---EBF becomes unfalsifiable
    \item \textbf{Lakatos violation:} Degenerative research program---no novel predictions
    \item \textbf{Cartwright violation:} Globalization of local effects---technocracy
\end{enumerate}

$\Rightarrow$ EBF degrades from science to policy rhetoric.

\subsubsection{The Position---Precisely Located}

The complementarity logic developed here is:
\begin{itemize}[noitemsep]
    \item \textbf{Not:} A new universal theory
    \item \textbf{Not:} A replacement for AD
    \item \textbf{But:} A methodological instrument for local marking of deviations under explicit retention of the null model
\end{itemize}

This position is:
\begin{itemize}[noitemsep]
    \item Popper-compatible (falsificationist)
    \item Lakatos-robust (progressive, not degenerative)
    \item Cartwright-realistic (local, not global)
    \item Fehr-compatible (exclusion, not integration)
\end{itemize}

\subsubsection{Final Synthesis: The Complete Position}

\begin{tcolorbox}[colback=blue!5!white,colframe=blue!75!black,title=\textbf{The Complete Philosophy of Science Position}]

\textbf{1. Three distinct levels exist:}
\begin{itemize}[noitemsep]
    \item Theory (trans-empirical, structures possibility space)
    \item Evidence (context-bound, locates deviations)
    \item Model choice (normative, decides application)
\end{itemize}

\textbf{2. Three foundational traditions ground this distinction:}
\begin{itemize}[noitemsep]
    \item Popper: Theories cannot be derived from data
    \item Lakatos: Hard cores require progressive protective belts
    \item Cartwright: Local causality does not globalize
\end{itemize}

\textbf{3. Category errors occur when levels are conflated:}
\begin{itemize}[noitemsep]
    \item Reducing theory to evidence (positivist error)
    \item Treating model choice as discovery (scientism error)
    \item Using evidence to ``replace'' theory (EBF's existential risk)
\end{itemize}

\textbf{4. The canonical relationship:}
\begin{quote}
\textit{Theories structure the space of possibilities.\\
Evidence localizes deviations.\\
EBF must not confuse the two.}
\end{quote}

\textbf{5. EBF's legitimate role:}
\begin{itemize}[noitemsep]
    \item Measuring $\gamma$ values (evidence)
    \item Using theoretical null models (presupposing theory)
    \item Informing model choice (supporting Level 3 decisions)
    \item \textit{Never} replacing theory (respecting Level 1 autonomy)
\end{itemize}

\end{tcolorbox}

\subsubsection{The Definitive Canonical Sentences}

\begin{tcolorbox}[colback=red!5!white,colframe=red!75!black,title=\textbf{Canonical Formulation (English)}]
\begin{center}
\Large
\textit{``Evidence localizes failure; theory structures possibility.\\
Confusing the two dissolves science.''}
\end{center}
\end{tcolorbox}

\begin{tcolorbox}[colback=red!5!white,colframe=red!75!black,title=\textbf{Kanonische Formulierung (Deutsch)}]
\begin{center}
\Large
\textit{``Empirie lokalisiert das Scheitern von Theorie---\\
sie ersetzt sie nicht.''}
\end{center}
\end{tcolorbox}

This is the philosophy of science essence of this appendix. Everything else is application of this principle.


% ============================================================================
% PRACTICAL IMPLICATIONS FOR COLLABORATION
% ============================================================================
\section{Practical Implications: Working with Fehr-Tradition Researchers}
\label{sec:practical}

This section translates the methodological principles into concrete guidance for academic collaboration.

\subsection{The Core Message}

\begin{tcolorbox}[colback=blue!5!white,colframe=blue!75!black,title=\textbf{Collaboration Condition}]
Collaboration with Ernst Fehr (or researchers in his tradition) is possible---and fruitful---\textbf{if and only if} it is clear that complementarity is a \textbf{bounding instrument}, not an integration promise.

Anything else will fail.
\end{tcolorbox}

\subsection{How Fehr Reads Your Work}

Fehr does not read charitably; he reads \textbf{disciplinedly}.

He always asks the same silent questions:
\begin{enumerate}
    \item What is the null hypothesis?
    \item What is explicitly excluded?
    \item Where can the model fail?
    \item What do I learn if it fails?
\end{enumerate}

If these four points are clear, he is open.\\
If not, he disengages internally.

\subsection{How to Frame Your Position}

\subsubsection{Dealbreakers (Never Say This)}

\begin{tcolorbox}[colback=red!5!white,colframe=red!75!black,title=\textbf{Phrases That Trigger Rejection}]
\begin{itemize}[noitemsep]
    \item ``We integrate all models.''
    \item ``Complementarity explains complex systems.''
    \item ``EBF replaces equilibrium theory.''
    \item ``Everything is connected to everything.''
\end{itemize}
These trigger exactly his red line.
\end{tcolorbox}

\subsubsection{Fehr-Compatible Framing (Say This)}

\begin{tcolorbox}[colback=green!5!white,colframe=green!75!black,title=\textbf{Phrases That Signal Methodological Rigor}]
\begin{itemize}[noitemsep]
    \item ``We begin with a restrictive null model.''
    \item ``Complementarity is a risky deviation hypothesis.''
    \item ``We show where $\gamma \approx 0$---and where not.''
    \item ``We explicitly exclude alternative mechanisms.''
\end{itemize}
This speaks directly his language.
\end{tcolorbox}

\subsection{The Role Fehr Would Take in Collaboration}

Important to understand: Fehr does \textbf{not} want to be:
\begin{itemize}[noitemsep]
    \item The visionary
    \item The integrator
    \item The systems theorist
\end{itemize}

He wants to be: \textbf{the methodological guardian of the discipline.}

In a good collaboration, he would be the one who asks:
\begin{itemize}[noitemsep]
    \item ``Why not simpler?''
    \item ``Where is the exclusion?''
    \item ``Is this really identified?''
\end{itemize}

This is not resistance---this is his contribution.

\subsection{How to Use Complementarity in a Fehr-Compatible Way}

The decisive reformulation:

\begin{center}
\textbf{Not:} ``Complementarity is everywhere.''\\[0.5em]
\textbf{But:} ``Complementarity is rare---and precisely therefore interesting.''
\end{center}

\textbf{Fehr accepts:}
\begin{itemize}[noitemsep]
    \item Rare, clearly identified interactions
    \item Strong null hypotheses
    \item Clean experimental designs
\end{itemize}

\textbf{Fehr rejects:}
\begin{itemize}[noitemsep]
    \item Diffuse systems rhetoric
    \item Integration ambitions
    \item Open meta-models
\end{itemize}

\subsection{Example: A Joint Project Structure}

\textbf{Title (Fehr-compatible):}
\begin{quote}
\textit{``When Complementarity Is Absent: Experimental Evidence from a No-Interaction Benchmark''}
\end{quote}

\textbf{Design:}
\begin{enumerate}
    \item Arrow-Debreu-style null model
    \item Strict assumption: $\gamma = 0$
    \item One targeted violation of this assumption
    \item Clear tests
    \item Acceptance of null findings
\end{enumerate}

Fehr loves null findings---if they are clean.

\subsection{What to Expect from Fehr}

Expect:
\begin{itemize}[noitemsep]
    \item Hard criticism
    \item Rejection of excessive complexity
    \item Demands for simplification
    \item Rejection of ``nice stories''
\end{itemize}

If you read this as an attack, the collaboration is over.\\
If you read it as quality assurance, it becomes very strong.

\subsection{The Unspoken Trust Question}

Fehr silently asks himself:

\begin{quote}
\textit{Does this person want to \textbf{explain} complexity---\\
or take \textbf{responsibility} for bounding it?}
\end{quote}

If your (implicit) answer is: ``I want to bound it,'' cooperation is possible.

\subsection{Key Sentences for Communication}

\begin{tcolorbox}[colback=gray!10!white,colframe=gray!75!black,title=\textbf{What You Could Say to Fehr}]

\textbf{Version 1:}\\
``I use complementarity not to integrate models, but to play them against each other and exclude them.''

\vspace{0.5em}
\textbf{Version 2 (stronger):}\\
``My goal is to show precisely when we do \textit{not} need complementarity.''

\end{tcolorbox}

This is exactly Fehr's point.

\subsection{The Realistic Summary}

\begin{enumerate}
    \item You must present your theory \textbf{narrower} than it is.
    \item You may keep the grand vision internally.
    \item In collaboration: \textbf{Discipline before synthesis.}
\end{enumerate}

If you accept this, Fehr becomes not an opponent but your strongest methodological ally.

\subsection{Final Sentence}

\begin{tcolorbox}[colback=blue!10!white,colframe=blue!75!black]
\begin{center}
\large
\textit{Ernst Fehr is not an opponent of complementarity---\\
he is the opponent of its \textbf{convenience}.}
\end{center}
\end{tcolorbox}


% ============================================================================
% KEY RESULTS
% ============================================================================
\section{Key Results}
\label{sec:results}

This section consolidates the main theoretical and methodological results of this appendix.

\subsection{Result 1: The Integration-Falsifiability Trade-off}

\begin{tcolorbox}[colback=gray!10!white,colframe=gray!75!black]
\textbf{Theorem (Informal):} As model integration capacity increases, falsifiability decreases---unless explicit exclusion commitments are made.

\textbf{Corollary:} A framework that can integrate all models explains nothing.
\end{tcolorbox}

\subsection{Result 2: The Exclusion Principle}
\label{sec:frm-exclusion-principle}

\begin{tcolorbox}[colback=gray!10!white,colframe=gray!75!black]
\textbf{Definition:} The Exclusion Principle states that $\gamma = 0$ is the null hypothesis for all complementarity claims. Scientific value comes from identifying where $\gamma \neq 0$, not from assuming non-zero interactions.

\textbf{Implication:} EBF gains falsifiability by committing to specific $\gamma = 0$ predictions.
\end{tcolorbox}

% =============================================================================
% FORMAL AXIOM: EXC-1 (Additivity Default)
% =============================================================================

\begin{axiom}[EXC-1: Additivity Default Principle]
\label{ax:exc-1}
For any aggregation of factors $\{f_1, \ldots, f_n\}$ into an outcome $Y$:

\begin{equation}
\boxed{
H_0: \quad Y = Y_0 + \sum_{j=1}^{n} \delta_j \cdot f_j \quad \textbf{(ADDITIVE = DEFAULT)}
}
\label{eq:exc-1-additive}
\end{equation}

The multiplicative alternative requires explicit justification:
\begin{equation}
H_1: \quad Y = Y_0 \cdot \prod_{j=1}^{n} f_j \quad \text{(only if Veto-Logic applies)}
\label{eq:exc-1-multiplicative}
\end{equation}
\end{axiom}

\begin{axiom}[EXC-2: Veto-Logic Criterion]
\label{ax:exc-2}
A multiplicative formulation is justified for factor $f_j$ if and only if the \textbf{Veto-Logic} holds:
\begin{equation}
\boxed{
f_j = 0 \quad \Rightarrow \quad Y = 0 \quad \text{(necessary condition for outcome)}
}
\label{eq:exc-2-veto}
\end{equation}

\textbf{Interpretation:} Factor $f_j$ is a ``gatekeeper''---without it, the outcome is impossible, not merely reduced.
\end{axiom}

\begin{axiom}[EXC-3: Hybrid Aggregation Rule]
\label{ax:exc-3}
When some factors satisfy Veto-Logic and others do not, use a \textbf{hybrid formulation}:
\begin{equation}
\boxed{
Y = Y_0 \cdot \prod_{j \in \mathcal{V}} f_j + \sum_{k \notin \mathcal{V}} \delta_k \cdot f_k
}
\label{eq:exc-3-hybrid}
\end{equation}
where $\mathcal{V} = \{j : f_j \text{ satisfies Veto-Logic}\}$.
\end{axiom}

\begin{axiom}[EXC-4: Documentation Requirement]
\label{ax:exc-4}
Any multiplicative formula in EBF \textbf{must} include a justification box:
\begin{enumerate}[noitemsep]
    \item \textbf{Veto-Logic Test:} Which factors satisfy $f_j = 0 \Rightarrow Y = 0$?
    \item \textbf{Empirical Evidence:} Citation showing interaction effect
    \item \textbf{Theoretical Mechanism:} Why do these factors interact (not just co-occur)?
    \item \textbf{Counterfactual:} What would the additive result be?
\end{enumerate}
Reference: Axiom EXC-1 (\ref{ax:exc-1})
\end{axiom}

\begin{axiom}[EXC-5: Exhaustive Complementarity Justifications (Whitelist)]
\label{ax:exc-5}
A multiplicative formulation is justified \textbf{if and only if} it satisfies \textbf{at least one} of the following ten criteria. Any justification not on this list defaults to additive (EXC-1).

\begin{center}
\renewcommand{\arraystretch}{1.3}
\small
\begin{tabular}{@{}clp{5.5cm}p{4cm}@{}}
\toprule
\textbf{\#} & \textbf{Code} & \textbf{Criterion} & \textbf{Test} \\
\midrule
\multicolumn{4}{l}{\textit{Veto \& Reversal (structural necessity):}} \\
1 & \textbf{V1} & \textbf{Veto-Logic}: $f_j = 0 \Rightarrow Y = 0$ & Factor=0 makes outcome impossible? \\
2 & \textbf{V2} & \textbf{Reversal-Logic}: $f_j < 0 \Rightarrow \text{sign}(Y)$ flips & Factor can reverse outcome? \\
\midrule
\multicolumn{4}{l}{\textit{Scaling \& Bounding (mathematical structure):}} \\
3 & \textbf{S1} & \textbf{Percentage-Scaling}: $f_j \in [1-\varepsilon, 1+\varepsilon]$, $\varepsilon \leq 0.5$ & Factor centered around 1? \\
4 & \textbf{S2} & \textbf{Bounded Output}: $Y = g(\prod f_j)$ with $g$ bounding & Product inside sigmoid/logit/etc.? \\
\midrule
\multicolumn{4}{l}{\textit{Physical \& Probabilistic (domain constraints):}} \\
5 & \textbf{P1} & \textbf{Dimensional Necessity}: Units require $\times$ & Does addition violate dimensions? \\
6 & \textbf{P2} & \textbf{Conditional Probability}: $P(A \cap B) = P(A) \cdot P(B|A)$ & Probability chain structure? \\
\midrule
\multicolumn{4}{l}{\textit{Gates \& Constraints (discrete mechanisms):}} \\
7 & \textbf{G1} & \textbf{Binary Gate}: $f_j \in \{0, 1\}$ & On/off switch mechanism? \\
8 & \textbf{C1} & \textbf{Capacity Constraint}: $f_j = \min(1, \text{cap}/\text{demand})$ & Exhaustible capacity? \\
\midrule
\multicolumn{4}{l}{\textit{Validated (empirical \& theoretical):}} \\
9 & \textbf{E1} & \textbf{Empirically Validated}: $H_0: \gamma=0$ rejected at $p<0.05$ & Additive statistically rejected? \\
10 & \textbf{T1} & \textbf{Theoretically Derived}: Axiom $\Rightarrow$ multiplicative & Follows from established theory? \\
\bottomrule
\end{tabular}
\end{center}

\textbf{Exclusion Rule:} Any justification \textit{not} matching V1--T1 is \textbf{inadmissible}. This includes but is not limited to: ``mathematically elegant,'' ``other papers do it,'' ``better empirical fit,'' ``intuitive,'' ``factors interact'' (without specifying how), ``conservative estimate,'' ``more general,'' ``captures nonlinearity'' (vague), ``factors are correlated,'' ``synergy effects'' (vague), ``diminishing returns'' (use concave function instead), ``cross-partial $\neq 0$'' (use additive interaction term).

\textbf{Default:} If no V1--T1 criterion is satisfied $\Rightarrow$ use additive (EXC-1).
\end{axiom}

\begin{table}[h!]
\centering
\caption{Exclusion Principle: Decision Matrix}
\label{tab:exc-decision-matrix}
\renewcommand{\arraystretch}{1.3}
\begin{tabular}{@{}lll@{}}
\toprule
\textbf{Question} & \textbf{Answer} & \textbf{Formulation} \\
\midrule
Does $f_j = 0 \Rightarrow Y = 0$? & Yes for all $j$ & Fully multiplicative \\
Does $f_j = 0 \Rightarrow Y = 0$? & Yes for some $j$ & Hybrid (EXC-3) \\
Does $f_j = 0 \Rightarrow Y = 0$? & No for all $j$ & Fully additive (default) \\
\bottomrule
\end{tabular}
\end{table}

\begin{example}[K* Formula: Exclusion Principle Application]
\label{ex:kstar-exclusion}
For the optimal portfolio size $K^*$ with factors Budget ($f_B$), Complexity ($f_C$), Stage ($f_S$), Crowding ($f_\gamma$), Scope ($f_L$):

\textbf{Veto-Logic Test:}
\begin{itemize}[noitemsep]
    \item $f_B = 0$ (no budget) $\Rightarrow K^* = 0$? \textbf{Yes} $\checkmark$ (no interventions possible)
    \item $f_C = 0$ (infinite complexity) $\Rightarrow K^* = 0$? \textbf{No} (fewer, not zero)
    \item $f_S = 0$ (wrong stage) $\Rightarrow K^* = 0$? \textbf{No} (less effective, not impossible)
    \item $f_\gamma = 0$ (severe crowding) $\Rightarrow K^* = 0$? \textbf{No} (caution, not prohibition)
    \item $f_L = 0$ (scope mismatch) $\Rightarrow K^* = 0$? \textbf{No} (adjustment, not veto)
\end{itemize}

\textbf{Result:} Only Budget satisfies Veto-Logic. Apply Hybrid Rule (EXC-3):
\begin{equation}
K^* = \underbrace{K_{\text{base}} \cdot f_B}_{\text{Veto (Budget)}} + \underbrace{\delta_C + \delta_S + \delta_\gamma + \delta_L}_{\text{Additive (Default)}}
\end{equation}
\end{example}

% =============================================================================
% CAUSAL FOUNDATIONS OF COMPLEMENTARITY (EXC-6)
% =============================================================================
\subsection{Causal Foundations: From Structural to Identified Complementarity}
\label{sec:causal-complementarity}

The EXC-5 whitelist provides \textit{structural} justifications for multiplicative formulations. This section establishes the \textit{causal} foundations that underpin empirical validation (E1) and connect structural complementarity to identifiable causal effects.

\begin{axiom}[EXC-6: Causal Identification Requirement]
\label{ax:exc-6}
A claim of \textbf{causal complementarity} requires both:
\begin{enumerate}[noitemsep]
    \item \textbf{Structural condition:} $\partial^2 f(x,z)/\partial x \partial z > 0$ (or equivalent)
    \item \textbf{Identification condition:} One of RCT, IV, or DiD (specified below)
\end{enumerate}
Without identification, any observed interaction is merely \textit{associative}, not \textit{causal}.
\end{axiom}

\subsubsection{Structural Model}

Let $X, Z$ be inputs and $Y$ the outcome. The true causal model is:
\begin{equation}
Y = f(X, Z, U)
\end{equation}
with unobserved disturbance $U$.

\begin{definition}[Technical Complementarity]
\label{def:technical-complementarity}
$X$ and $Z$ exhibit \textbf{technical complementarity} if:
\begin{equation}
\frac{\partial^2 f(x,z,u)}{\partial x\,\partial z} > 0 \quad \forall (x,z,u)
\end{equation}
This is a purely structural property of $f$, independent of identification.
\end{definition}

\subsubsection{Potential Outcomes Formalism}

Define potential outcomes:
\begin{equation}
Y(x,z) := \text{Outcome when } X=x \text{ and } Z=z
\end{equation}

\begin{definition}[Causal Complementarity]
\label{def:causal-complementarity}
$X$ and $Z$ are \textbf{causally complementary} if for all $x' > x$ and $z' > z$:
\begin{equation}
\underbrace{\big[ Y(x', z') - Y(x, z') \big]}_{\Delta_X(z')} > \underbrace{\big[ Y(x', z) - Y(x, z) \big]}_{\Delta_X(z)}
\end{equation}
Equivalently: $\Delta_X(z') > \Delta_X(z)$ --- the causal effect of $X$ is increasing in $Z$.
\end{definition}

For continuous variables:
\begin{equation}
\frac{\partial}{\partial z}\left(\frac{\partial \mathbb{E}[Y(x,z)]}{\partial x}\right) > 0
\end{equation}

\subsubsection{Pearl's do-Calculus Formulation}

The most precise causal definition uses intervention notation:
\begin{equation}
\boxed{\frac{\partial}{\partial z}\left(\frac{\partial \mathbb{E}[Y \mid do(X=x), do(Z=z)]}{\partial x}\right) > 0}
\end{equation}

\begin{tcolorbox}[colback=blue!5!white,colframe=blue!75!black,title=Causal Complementarity Definition]
\textbf{Complementarity is causal} when the effect of an intervention on $X$ depends on the intervention level of $Z$.
\end{tcolorbox}

\subsubsection{Increasing Differences and Supermodularity}

\begin{definition}[Increasing Differences]
\label{def:increasing-differences}
For ordered sets $\mathcal{X}, \mathcal{Z}$, function $m(x,z)$ has \textbf{increasing differences} (ID) in $(x,z)$ if:
\begin{equation}
m(x',z) - m(x,z) \text{ is (weakly) increasing in } z \quad \forall x' > x
\tag{ID}
\end{equation}
This is exactly the potential-outcomes definition of causal complementarity.
\end{definition}

\begin{proposition}[ID $\Leftrightarrow$ Supermodularity (Discrete)]
\label{prop:id-supermodular}
For all $x' > x$ and $z' > z$:
\begin{equation}
m \text{ has ID} \quad \Longleftrightarrow \quad m(x',z') + m(x,z) \geq m(x',z) + m(x,z')
\tag{SM}
\end{equation}
\end{proposition}

\begin{proof}[Proof Sketch]
(ID) means $m(x',z') - m(x,z') \geq m(x',z) - m(x,z)$.
Rearranging: $m(x',z') + m(x,z) \geq m(x',z) + m(x,z')$, which is (SM).
The reverse follows by back-substitution. \qed
\end{proof}

\begin{proposition}[ID $\Rightarrow$ Cross-Derivative $\geq 0$ (Continuous)]
\label{prop:id-crossderiv}
If $m$ is twice continuously differentiable, then:
\begin{equation}
m \text{ has ID} \quad \Longrightarrow \quad \frac{\partial^2 m(x,z)}{\partial x\,\partial z} \geq 0 \quad \forall(x,z)
\tag{CA}
\end{equation}
\end{proposition}

\begin{proof}[Proof Sketch]
ID means $\partial m/\partial x$ is increasing in $z$.
``Increasing in $z$'' means $\partial(\partial m/\partial x)/\partial z \geq 0$, i.e., $\partial^2 m/\partial x\partial z \geq 0$. \qed
\end{proof}

\subsubsection{Identification Conditions}

\begin{lemma}[Identification under Joint Ignorability]
\label{lem:joint-ignorability}
Under the assumption:
\begin{equation}
(X,Z) \perp\!\!\!\perp Y(x,z) \quad \forall(x,z)
\tag{A0}
\end{equation}
causal complementarity is identified by:
\begin{equation}
m(x,z) = \mathbb{E}[Y \mid X=x, Z=z]
\end{equation}
and is testable via supermodularity inequalities in observed cell means.
\end{lemma}

For binary $X, Z \in \{0,1\}$, with $\mu_{xz} := \mathbb{E}[Y \mid X=x, Z=z]$:
\begin{equation}
\boxed{\Delta := (\mu_{11} - \mu_{01}) - (\mu_{10} - \mu_{00}) \geq 0}
\end{equation}

\subsubsection{Difference-in-Differences Identification}

For panel data with $Y_{it}$, time $t \in \{0,1\}$, treatment $X_i \in \{0,1\}$, and moderator $Z_i \in \{0,1\}$:

\begin{definition}[Causal DiD-Complementarity]
\label{def:did-complementarity}
\begin{multline}
\Big(\mathbb{E}[Y_{i1}(1,1) - Y_{i0}(1,1)] - \mathbb{E}[Y_{i1}(0,1) - Y_{i0}(0,1)]\Big) \\
- \Big(\mathbb{E}[Y_{i1}(1,0) - Y_{i0}(1,0)] - \mathbb{E}[Y_{i1}(0,0) - Y_{i0}(0,0)]\Big) \geq 0
\tag{DiD-ID}
\end{multline}
Interpretation: ``The DiD effect of $X$ is larger when $Z=1$.''
\end{definition}

\begin{assumption}[Stratified Parallel Trends]
\label{ass:parallel-trends}
For each $z$:
\begin{equation}
\mathbb{E}[Y_{i1}(0,z) - Y_{i0}(0,z) \mid X=1, Z=z] = \mathbb{E}[Y_{i1}(0,z) - Y_{i0}(0,z) \mid X=0, Z=z]
\tag{PT(z)}
\end{equation}
\end{assumption}

Under PT(z), the triple-difference regression identifies causal complementarity:
\begin{equation}
Y_{it} = \alpha_i + \lambda_t + \theta Z_i + \beta (X_i \times \text{Post}_t) + \phi (Z_i \times \text{Post}_t) + \delta (X_i \times Z_i \times \text{Post}_t) + \varepsilon_{it}
\end{equation}
where $\delta$ is the \textbf{causally identified complementarity coefficient}.

\subsubsection{Topkis Theorem: Policy Implications}

\begin{theorem}[Topkis, Monotone Comparative Statics]
\label{thm:topkis}
Let $x$ be a decision (e.g., investment intensity), $z$ a parameter (e.g., technology level).
If:
\begin{enumerate}[noitemsep]
    \item $\mathcal{X}$ is a lattice (interval or ordered discrete set)
    \item $m(x,z)$ is supermodular / has increasing differences
\end{enumerate}
then any selection $x^*(z) \in \arg\max_x m(x,z)$ is \textbf{monotonically increasing} in $z$.
\end{theorem}

\begin{proof}[Proof Idea]
Supermodularity means ``high $x$'' and ``high $z$'' are mutually attractive. If higher $x$ is optimal at $z$, it cannot become optimal to switch to lower $x$ at $z' > z$, because the relative advantage of higher $x$ increases with $z$. \qed
\end{proof}

\textbf{Implication for EBF:} If $m(x,z) = \mathbb{E}[Y \mid do(X=x), do(Z=z)]$ is supermodular, then:
\begin{itemize}[noitemsep]
    \item \textbf{Policy complementarity:} At higher $z$, one should (causally) choose higher $x$
    \item \textbf{Treatment complementarity:} The treatment effect of $X$ is larger when $Z$ is higher
\end{itemize}

\subsubsection{Empirical Test Summary}

\begin{tcolorbox}[colback=green!5!white,colframe=green!75!black,title=Causal Complementarity: Test Formulas]
\textbf{Binary (2×2):}
\begin{equation}
\Delta = \mu_{11} + \mu_{00} - \mu_{10} - \mu_{01} \geq 0
\end{equation}
with $\mu_{xz} = \mathbb{E}[Y \mid do(X=x), do(Z=z)]$.

\textbf{Continuous (local):}
\begin{equation}
\frac{\partial^2}{\partial x \partial z}\, \mathbb{E}[Y \mid do(X=x), do(Z=z)] \geq 0
\end{equation}

\textbf{DiD (with time):}
\begin{equation}
\delta \text{ in triple-interaction term } (X \times Z \times \text{Post}) \geq 0
\end{equation}
\end{tcolorbox}

\subsubsection{Connection to EXC-5}

The causal framework (EXC-6) enriches the EXC-5 whitelist:
\begin{itemize}[noitemsep]
    \item \textbf{V1, V2, S1, S2, G1, C1:} Structural mechanisms (directly verifiable)
    \item \textbf{P1, P2:} Dimensional/probabilistic necessity (domain constraints)
    \item \textbf{E1:} Requires causal identification per EXC-6 (RCT, IV, or DiD)
    \item \textbf{T1:} Theoretical derivation implies causal structure
\end{itemize}

\begin{table}[h!]
\centering
\caption{EXC-5 Codes and Identification Requirements}
\label{tab:exc5-identification}
\renewcommand{\arraystretch}{1.2}
\small
\begin{tabular}{@{}llll@{}}
\toprule
\textbf{Code} & \textbf{Type} & \textbf{Identification} & \textbf{Test} \\
\midrule
V1, V2 & Structural & Direct mechanism & $f=0 \Rightarrow Y=0$? \\
S1, S2 & Mathematical & Functional form & Bounded? Centered at 1? \\
P1, P2 & Physical/Prob. & Domain constraint & Units? Probability chain? \\
G1, C1 & Discrete & Mechanism & Binary? Capacity limit? \\
E1 & Empirical & \textbf{EXC-6 required} & RCT/IV/DiD \\
T1 & Theoretical & Theory $\Rightarrow$ structure & Axiom derivation \\
\bottomrule
\end{tabular}
\end{table}

\subsubsection{Quick Reference: Multiplicative vs. Additive Decision}

\begin{tcolorbox}[colback=yellow!5!white,colframe=orange!75!black,title=Decision Flowchart: When to Use Multiplicative]
\textbf{Default Rule:} If none of the following apply $\Rightarrow$ \textbf{ADDITIVE} (EXC-1).

\begin{center}
\renewcommand{\arraystretch}{1.4}
\small
\begin{tabular}{@{}clcp{5cm}@{}}
\toprule
\textbf{\#} & \textbf{Question} & \textbf{If YES} & \textbf{Test} \\
\midrule
1 & Does $f=0 \Rightarrow Y=0$? (impossible, not just reduced) & \textbf{V1} $\times$ & Budget=0 $\Rightarrow$ K*=0? \\
2 & Can $f<0$ reverse the outcome (backfire)? & \textbf{V2} $\times$ & $\sigma_s < 0 \Rightarrow E < 0$? \\
3 & Are ALL factors in $[1-\varepsilon, 1+\varepsilon]$, $\varepsilon \leq 0.5$? & \textbf{S1} $\times$ & $g(\varphi) \in [0.8, 1.3]$? \\
4 & Is the product INSIDE a bounding function? & \textbf{S2} $\times$ & $\sigma(\text{WEC} \cdot \alpha \cdot \beta)$? \\
5 & Do UNITS require multiplication? & \textbf{P1} $\times$ & Area = L $\times$ W? \\
6 & Is it a PROBABILITY CHAIN $P(A) \cdot P(B|A)$? & \textbf{P2} $\times$ & Sequential stages? \\
7 & Is the factor BINARY $\{0, 1\}$? & \textbf{G1} $\times$ & On/off switch? \\
8 & Is there an EXHAUSTIBLE CAPACITY? & \textbf{C1} $\times$ & Attention $\times$ Salience? \\
9 & Did RCT/IV/DiD REJECT additive ($p<0.05$)? & \textbf{E1} $\times$ & Causal identification (EXC-6)? \\
10 & Does THEORY axiomatically require $\times$? & \textbf{T1} $\times$ & Fehr-Schmidt $\Rightarrow \alpha \cdot \beta$? \\
\midrule
\multicolumn{2}{l}{\textbf{None of the above?}} & \textbf{ADD} & $Y = Y_0 + \sum_j \delta_j \cdot f_j$ \\
\bottomrule
\end{tabular}
\end{center}

\textbf{Hybrid Rule (EXC-3):} If SOME factors satisfy V1/V2/G1 and others don't:
\begin{equation}
Y = \underbrace{Y_0 \cdot \prod_{j \in \text{Veto}} f_j}_{\text{Multiplicative}} + \underbrace{\sum_{k \notin \text{Veto}} \delta_k \cdot f_k}_{\text{Additive}}
\end{equation}
\end{tcolorbox}

\begin{tcolorbox}[colback=red!5!white,colframe=red!75!black,title=Common Mistakes (Auto-Reject)]
The following arguments are \textbf{NEVER} valid justifications for multiplicative:
\begin{itemize}[noitemsep]
    \item ``Looks more elegant'' (X1) $\rightarrow$ Aesthetics $\neq$ accuracy
    \item ``Famous paper X does it'' (X4) $\rightarrow$ Authority $\neq$ evidence
    \item ``Factors are correlated'' (Y2) $\rightarrow$ Correlation $\neq$ complementarity
    \item ``Diminishing returns'' (Y3) $\rightarrow$ Use concave $f(x) + g(y)$
    \item ``Synergy'' (Y5) $\rightarrow$ Vague; specify V1/V2/S1/etc.
    \item ``Factor is important'' (Y1) $\rightarrow$ Importance $\neq$ veto
\end{itemize}
\textbf{Rule:} If your justification isn't V1--T1, it defaults to additive.
\end{tcolorbox}

\subsubsection{Worked Examples: EXC-5 Analysis in Practice}

The following ten examples demonstrate how to apply the EXC-5 decision process systematically.

\begin{example}[Donation Willingness]
\label{ex:exc-donation}
\textbf{Proposed:} $D = \text{Income} \cdot \text{Empathy} \cdot \text{Visibility}$

\textbf{Veto Analysis:}
\begin{itemize}[noitemsep]
    \item Income = 0 $\Rightarrow$ D = 0? \textbf{Yes} (no money = no donation possible) $\checkmark$ V1
    \item Empathy = 0 $\Rightarrow$ D = 0? \textbf{No} (social motives, tax benefits still work)
    \item Visibility = 0 $\Rightarrow$ D = 0? \textbf{No} (anonymous donations exist)
\end{itemize}

\textbf{Result:} Only Income satisfies V1. $\Rightarrow$ \textbf{Hybrid (EXC-3):}
\begin{equation}
D = D_{\text{base}} \cdot f_{\text{Income}} + \delta_{\text{Empathy}} + \delta_{\text{Visibility}}
\end{equation}
\end{example}

\begin{example}[Workplace Productivity]
\label{ex:exc-productivity}
\textbf{Proposed:} $P = \text{Motivation} \cdot \text{Ability} \cdot \text{Resources}$

\textbf{Veto Analysis:}
\begin{itemize}[noitemsep]
    \item Motivation = 0 $\Rightarrow$ P = 0? \textbf{No} (fear of termination drives minimum effort)
    \item Ability = 0 $\Rightarrow$ P = 0? \textbf{Yes} (no skill = no production) $\checkmark$ V1
    \item Resources = 0 $\Rightarrow$ P = 0? \textbf{Yes} (no tools = no work possible) $\checkmark$ V1
\end{itemize}

\textbf{Result:} Ability and Resources satisfy V1. $\Rightarrow$ \textbf{Hybrid:}
\begin{equation}
P = P_{\text{base}} \cdot f_{\text{Ability}} \cdot f_{\text{Resources}} + \delta_{\text{Motivation}}
\end{equation}
\end{example}

\begin{example}[Purchase Probability]
\label{ex:exc-purchase}
\textbf{Proposed:} $P_{\text{buy}} = \text{Awareness} \cdot \text{Interest} \cdot \text{Budget}$

\textbf{Veto Analysis:}
\begin{itemize}[noitemsep]
    \item Awareness = 0 $\Rightarrow$ P = 0? \textbf{Yes} (unknown product cannot be bought) $\checkmark$ V1
    \item Interest = 0 $\Rightarrow$ P = 0? \textbf{No} (impulse purchases without genuine interest)
    \item Budget = 0 $\Rightarrow$ P = 0? \textbf{Yes} (no money = no purchase) $\checkmark$ V1
\end{itemize}

\textbf{Result:} Awareness and Budget satisfy V1. $\Rightarrow$ \textbf{Hybrid:}
\begin{equation}
P_{\text{buy}} = P_{\text{base}} \cdot f_A \cdot f_B + \delta_{\text{Interest}}
\end{equation}
\end{example}

\begin{example}[Team Performance --- Fully Additive]
\label{ex:exc-team}
\textbf{Proposed:} $T = \text{Skill}_1 \cdot \text{Skill}_2 \cdot \text{Communication}$

\textbf{Veto Analysis:}
\begin{itemize}[noitemsep]
    \item Skill$_1$ = 0 $\Rightarrow$ T = 0? \textbf{No} (team can compensate)
    \item Skill$_2$ = 0 $\Rightarrow$ T = 0? \textbf{No} (team can compensate)
    \item Communication = 0 $\Rightarrow$ T = 0? \textbf{No} (poor communication $\neq$ zero output)
\end{itemize}

\textbf{Result:} No factor satisfies V1. $\Rightarrow$ \textbf{Fully Additive (EXC-1):}
\begin{equation}
T = T_{\text{base}} + \delta_1 + \delta_2 + \delta_{\text{Comm}}
\end{equation}
\end{example}

\begin{example}[Medication Effectiveness --- Fully Multiplicative]
\label{ex:exc-medication}
\textbf{Proposed:} $W = \text{Dose} \cdot \text{Compliance} \cdot \text{Absorption}$

\textbf{Veto Analysis:}
\begin{itemize}[noitemsep]
    \item Dose = 0 $\Rightarrow$ W = 0? \textbf{Yes} (no dose = no effect) $\checkmark$ V1
    \item Compliance = 0 $\Rightarrow$ W = 0? \textbf{Yes} (not taken = no effect) $\checkmark$ V1
    \item Absorption = 0 $\Rightarrow$ W = 0? \textbf{Yes} (not absorbed = no effect) $\checkmark$ V1
\end{itemize}

\textbf{Result:} ALL factors satisfy V1. $\Rightarrow$ \textbf{Fully Multiplicative:}
\begin{equation}
W = W_{\text{max}} \cdot f_D \cdot f_C \cdot f_A
\end{equation}
\end{example}

\begin{example}[Conversion Rate --- S2 Bounded Output]
\label{ex:exc-conversion}
\textbf{Proposed:} $C = \sigma(\text{Reach} \cdot \text{Relevance} \cdot \text{Timing})$

\textbf{Analysis:}
\begin{itemize}[noitemsep]
    \item Product is \textit{inside} sigmoid function $\sigma(\cdot)$
    \item Reach = 0? $\Rightarrow \sigma(0) = 0.5$, not 0 (no veto)
    \item Output bounded to $(0, 1)$ regardless of input magnitudes
\end{itemize}

\textbf{Result:} S2 (Bounded Output) applies. $\Rightarrow$ \textbf{Multiplicative justified:}
\begin{equation}
C = \sigma(\text{Reach} \cdot \text{Relevance} \cdot \text{Timing})
\end{equation}
\textbf{Code: S2}
\end{example}

\begin{example}[Intervention with Backfire Risk --- V2 Reversal]
\label{ex:exc-backfire}
\textbf{Proposed:} $E = E_{\text{base}} \cdot \alpha_{\text{phase}} \cdot \sigma_{\text{segment}}$

\textbf{Analysis:}
\begin{itemize}[noitemsep]
    \item $\alpha_{\text{phase}} \in [0.3, 0.9]$ --- cannot be negative, just scaling
    \item $\sigma_{\text{segment}} \in [-0.5, 2.0]$ --- \textbf{CAN BE NEGATIVE} (reactance $\Rightarrow$ backfire)
\end{itemize}

\textbf{Result:} $\sigma_s < 0$ causes sign reversal. $\Rightarrow$ \textbf{V2 (Reversal) applies:}
\begin{equation}
E = E_{\text{base}} \cdot \alpha \cdot \sigma_s
\end{equation}
\textbf{Code: V2}
\end{example}

\begin{example}[Dynamic Portfolio --- S1 Percentage Scaling]
\label{ex:exc-dynamic}
\textbf{Proposed:} $K^*(t) = K^*_0 \cdot g(\varphi) \cdot h(\Psi) \cdot \ell(\Theta)$

\textbf{Analysis:}
\begin{itemize}[noitemsep]
    \item $g(\varphi) = 1 + 0.3 \cdot |\Delta\varphi| \in [1.0, 1.3]$ --- centered around 1
    \item $h(\Psi) \in \{1.0, 1.5\}$ --- centered around 1
    \item $\ell(\Theta) = 1 - 0.2 \cdot \text{converged} \in [0.8, 1.0]$ --- centered around 1
\end{itemize}

\textbf{Result:} All factors in $[1 \pm 0.5]$. $\Rightarrow$ \textbf{S1 (Percentage Scaling) applies:}
\begin{equation}
K^*(t) = K^*_0 \cdot g \cdot h \cdot \ell
\end{equation}
\textbf{Code: S1}
\end{example}

\begin{example}[Voter Turnout]
\label{ex:exc-turnout}
\textbf{Proposed:} $V = \text{Interest} \cdot \text{Accessibility} \cdot \text{Weather}$

\textbf{Veto Analysis:}
\begin{itemize}[noitemsep]
    \item Interest = 0 $\Rightarrow$ V = 0? \textbf{No} (civic duty even without interest)
    \item Accessibility = 0 $\Rightarrow$ V = 0? \textbf{Yes} (no access = no vote possible) $\checkmark$ V1
    \item Weather = 0 (bad) $\Rightarrow$ V = 0? \textbf{No} (reduces, doesn't prevent)
\end{itemize}

\textbf{Result:} Only Accessibility satisfies V1. $\Rightarrow$ \textbf{Hybrid:}
\begin{equation}
V = V_{\text{base}} \cdot f_{\text{Access}} + \delta_{\text{Interest}} + \delta_{\text{Weather}}
\end{equation}
\end{example}

\begin{example}[Learning Effect]
\label{ex:exc-learning}
\textbf{Proposed:} $L = \text{Practice} \cdot \text{Feedback} \cdot \text{Sleep}$

\textbf{Veto Analysis:}
\begin{itemize}[noitemsep]
    \item Practice = 0 $\Rightarrow$ L = 0? \textbf{Yes} (no practice = no learning) $\checkmark$ V1
    \item Feedback = 0 $\Rightarrow$ L = 0? \textbf{No} (slower learning without feedback, but possible)
    \item Sleep = 0 $\Rightarrow$ L = 0? \textbf{No} (sleep-deprived can still learn, inefficiently)
\end{itemize}

\textbf{Result:} Only Practice satisfies V1. $\Rightarrow$ \textbf{Hybrid:}
\begin{equation}
L = L_{\text{base}} \cdot f_{\text{Practice}} + \delta_{\text{Feedback}} + \delta_{\text{Sleep}}
\end{equation}
\end{example}

\begin{table}[h!]
\centering
\caption{Summary: 10 EXC-5 Worked Examples}
\label{tab:exc-examples-summary}
\renewcommand{\arraystretch}{1.2}
\small
\begin{tabular}{@{}clll@{}}
\toprule
\textbf{\#} & \textbf{Domain} & \textbf{Result} & \textbf{Code} \\
\midrule
1 & Donation & Hybrid & V1 (Income) \\
2 & Productivity & Hybrid & V1 (Ability, Resources) \\
3 & Purchase & Hybrid & V1 (Awareness, Budget) \\
4 & Team & \textbf{Additive} & --- \\
5 & Medication & Multiplicative & V1 (all 3) \\
6 & Conversion & Multiplicative & S2 (sigmoid) \\
7 & Intervention & Multiplicative & V2 ($\sigma_s < 0$) \\
8 & Portfolio & Multiplicative & S1 (around 1) \\
9 & Turnout & Hybrid & V1 (Accessibility) \\
10 & Learning & Hybrid & V1 (Practice) \\
\bottomrule
\end{tabular}
\end{table}

\subsection{Result 3: The Three-Level Distinction}

\begin{tcolorbox}[colback=gray!10!white,colframe=gray!75!black]
\textbf{Proposition:} Scientific knowledge operates at three distinct levels:
\begin{enumerate}[noitemsep]
    \item \textbf{Theory} (trans-empirical, structures possibility)
    \item \textbf{Evidence} (context-bound, localizes failure)
    \item \textbf{Model Choice} (normative, cannot be resolved by evidence alone)
\end{enumerate}

\textbf{Implication:} Confusing these levels dissolves scientific practice.
\end{tcolorbox}

\subsection{Result 4: Non-Integrability as Discipline}

\begin{tcolorbox}[colback=gray!10!white,colframe=gray!75!black]
\textbf{Theorem (Informal):} Not a single null model disciplines EBF, but the explicit acknowledgment of non-integrability between models.

\textbf{Method:} Work with multiple incompatible models; show EBF can rationalize each individually but not jointly.
\end{tcolorbox}

\subsection{Result 5: Additivity-Complementarity Equivalence}

\begin{tcolorbox}[colback=gray!10!white,colframe=gray!75!black]
\textbf{Proposition:} Additivity and complementarity are equally valid scientific perspectives with different epistemic functions:
\begin{itemize}[noitemsep]
    \item Additivity structures \textit{knowledge}
    \item Complementarity structures \textit{application}
\end{itemize}

\textbf{Practical Rule:} Start additive; switch to complementarity when additivity systematically fails.
\end{tcolorbox}

\subsection{Summary Table of Results}

\begin{table}[h!]
\centering
\renewcommand{\arraystretch}{1.3}
\begin{tabular}{clp{7cm}}
\toprule
\textbf{\#} & \textbf{Result} & \textbf{Implication for EBF} \\
\midrule
R1 & Integration-Falsifiability Trade-off & Commit to exclusions \\
R2 & Exclusion Principle & $\gamma = 0$ as null hypothesis \\
R3 & Three-Level Distinction & Don't confuse theory with evidence \\
R4 & Non-Integrability Discipline & Work with multiple incompatible models \\
R5 & Additivity-Complementarity Equivalence & Both valid; different roles \\
\bottomrule
\end{tabular}
\caption{Summary of key theoretical results}
\end{table}


% ============================================================================
% METHODOLOGICAL REFINEMENT: MULTIPLE MODELS STRATEGY
% ============================================================================
\section{Methodological Refinement: Multiple Models Instead of a Single Null}
\label{sec:multiple-models}

This section presents a crucial methodological refinement. Instead of explicitly referencing
Arrow-Debreu as \textit{the} null model, a more elegant and less confrontational strategy
works with \textbf{multiple clearly incompatible models}.

\subsection{Why Avoiding Explicit Arrow-Debreu Reference Is Often Preferable}

Arrow-Debreu is methodologically useful but rhetorically charged.

\textbf{Problems with explicit AD reference:}
\begin{itemize}[noitemsep]
    \item Triggers immediate paradigm conflict (``Neoclassical vs. Behavioral'')
    \item Provokes defensive reactions from both camps
    \item Leads into false debates about ``realism''
    \item Forces Fehr (and others) into a ``defender'' role they don't want
\end{itemize}

$\Rightarrow$ \textbf{This distracts from the actual point.}

The central claim is \textit{not}:

\begin{quote}
``Arrow-Debreu is wrong.''
\end{quote}

But rather:

\begin{quote}
\textbf{``EBF cannot replace model choice.''}
\end{quote}

This can be demonstrated more cleanly \textit{without} AD.

\subsection{The Better Strategy: Multiple Restrictive Models}

Instead of one explicit null model (AD), work with:

\begin{tcolorbox}[colback=green!5!white,colframe=green!75!black,title=\textbf{The Multiple Models Strategy}]
Present \textbf{multiple clearly incompatible} economic and psychological models
and show that EBF can rationalize each individually---but not jointly.

This is methodologically even \textit{stronger} than a single null model.
\end{tcolorbox}

\subsubsection{Example Model Palette}

\textbf{Economic Models:}
\begin{itemize}[noitemsep]
    \item Standard incentive model (monetary incentives)
    \item Social preferences (inequality aversion)
    \item Reciprocity model
    \item Norms/expectations model
\end{itemize}

\textbf{Psychological Models:}
\begin{itemize}[noitemsep]
    \item Intrinsic motivation
    \item Framing effects
    \item Self-image / Identity
    \item Cognitive limitations
\end{itemize}

$\Rightarrow$ These models are:
\begin{itemize}[noitemsep]
    \item Each individually plausible
    \item Each well-supported empirically
    \item But \textbf{not jointly integrable} without losing falsifiability
\end{itemize}

\subsection{What EBF Can Do (State Clearly)}

\textbf{EBF can:}
\begin{enumerate}[noitemsep]
    \item \textbf{Identify local effects}: ``In context $X$, incentive $A$ works''
    \item \textbf{Classify contexts}: ``Here norm effects appear dominant''
    \item \textbf{Measure interactions locally}: ``$A \times B$ is significant here / not there''
\end{enumerate}

$\Rightarrow$ \textbf{EBF excels at empirical mapping.}

\subsection{What EBF Cannot Do (Make Explicit)}

\textbf{EBF cannot:}
\begin{enumerate}[noitemsep]
    \item Decide between incompatible models
    \item Combine models without losing identification
    \item Say which model is fundamentally correct
    \item Replace theory
    \item Neutralize normative model choice
\end{enumerate}

$\Rightarrow$ \textbf{This is not a deficiency but a boundary.}

\subsection{The Central Point Without Arrow-Debreu}

\begin{tcolorbox}[colback=blue!5!white,colframe=blue!75!black,title=\textbf{The Formulation (AD-Free)}]
\begin{center}
\large
\textit{Evidence-based methods can rationalize multiple competing models in isolation.\\
Their strength lies in local empirical validation,\\
not in global model integration or selection.}
\end{center}

\vspace{0.5em}
Or in German:

\begin{center}
\textit{EBF kann zeigen, dass jedes der gängigen Modelle in bestimmten Kontexten funktioniert---\\
aber genau deshalb kann sie nicht entscheiden,\\
welche Modelle gemeinsam gelten dürfen.}
\end{center}
\end{tcolorbox}

\subsection{Where Complementarity Now Fits Cleanly}

With this framing, complementarity becomes \textit{not}:

\begin{quote}
$\times$ ``Everything is connected to everything''
\end{quote}

But rather:

\begin{quote}
$\checkmark$ A tool to show precisely \textbf{which model combinations are impermissible}
because they would destroy all falsification.
\end{quote}

Complementarity becomes:
\begin{itemize}[noitemsep]
    \item \textbf{Selective} (not universal)
    \item \textbf{Ex-ante bounded} (not post-hoc accommodating)
    \item \textbf{Model-separating} (not model-unifying)
\end{itemize}

\subsection{The Methodologically Clean Narrative (Step-by-Step)}

\begin{enumerate}
    \item \textbf{Present} multiple plausible models side by side
    \item \textbf{Show} EBF results that support each model individually
    \item \textbf{Demonstrate} that the same evidence supports multiple models
    \item \textbf{Show} that their combination is not identifiable
    \item \textbf{Conclude:}
    \begin{itemize}[noitemsep]
        \item EBF explains effects
        \item Theory constrains models
        \item Complementarity must not integrate
    \end{itemize}
\end{enumerate}

$\Rightarrow$ \textbf{This is unassailable---including for Fehr.}

\subsection{Canonical Formulations for Academic Use}

\begin{tcolorbox}[colback=green!5!white,colframe=green!75!black]
\begin{center}
\textit{``The contribution of evidence-based approaches is not to unify behavioral models,\\
but to reveal the limits of their joint applicability.''}
\end{center}
\end{tcolorbox}

Or more concisely:

\begin{tcolorbox}[colback=blue!5!white,colframe=blue!75!black]
\begin{center}
\large
\textit{``EBF tells us where models work---\\
not which models belong together.''}
\end{center}
\end{tcolorbox}

\subsection{What This Gains}

This strategy achieves:
\begin{itemize}[noitemsep]
    \item No paradigm wars
    \item No Arrow-Debreu front
    \item Full compatibility with: Fehr, Behavioral Economics, Policy
    \item Clear philosophy of science position
\end{itemize}

\subsection{The Refined Disciplining Principle}

\begin{tcolorbox}[colback=red!5!white,colframe=red!75!black,title=\textbf{The Key Insight (Memorize)}]
\begin{center}
\Large
\textit{Not a single null model disciplines EBF,\\
but the explicit acknowledgment of\\
\textbf{non-integrability between models}.}
\end{center}

\vspace{0.5em}
\normalsize
German:

\begin{center}
\textit{Nicht ein Nullmodell diszipliniert EBF,\\
sondern die explizite Anerkennung von\\
\textbf{Nicht-Integrierbarkeit zwischen Modellen}.}
\end{center}
\end{tcolorbox}

This is the refined methodological position that builds on Section~\ref{sec:taxonomy}
and makes the Exclusion Principle operational without requiring commitment to any
particular theoretical baseline.


% ============================================================================
% GLOSSARY
% ============================================================================
\section{Glossary}
\label{sec:glossary}

Key terms (see also Appendix GLS):

\begin{description}[style=nextline]
    \item[Integration Principle] Using $\gamma$ to explain why elements fit together; risks unfalsifiability
    \item[Exclusion Principle] Using $\gamma = 0$ as null hypothesis; maintains falsifiability
    \item[Exclusion Zone] Set of pairs $(a,b)$ for which framework commits to $\gamma_{ab} = 0$
    \item[Integrable Models] Models that can be combined under clear conditions (additive baseline + deviation, within-mechanism complementarity, temporal sequencing)
    \item[Non-Integrable Models] Models that cannot be combined without identification collapse, logical contradiction, or normative incoherence
    \item[Disciplinary Power] The capacity of a model to exclude observations; source of scientific value
    \item[Progressive Research Program] (Lakatos) Program that predicts novel facts and learns from failure
    \item[Degenerative Research Program] (Lakatos) Program that protects core by ad-hoc auxiliary hypotheses
    \item[Framework] Organizing structure for knowledge; not directly testable
    \item[Model] Specific instantiation with testable predictions
    \item[Trans-Empirical] Knowledge that is logically prior to observation; structures interpretation of data rather than being derived from it (e.g., theoretical frameworks)
    \item[Category Error] Treating one level of scientific knowledge as if it were another; e.g., reducing theory to evidence
    \item[Null Model] A theoretical construct defining what would happen if certain effects were absent; enables identification and measurement of deviations
    \item[Logical Positivism] Failed philosophical program claiming all meaningful statements are either analytic or empirically verifiable; relevant as historical parallel to EBF risks
    \item[Model Choice] Level 3 scientific knowledge; the normative decision about which theory to apply in which context; not reducible to evidence alone
    \item[Hard Core] (Lakatos) The fundamental commitments of a research program that are never directly abandoned; for economics, Arrow-Debreu serves this function
    \item[Protective Belt] (Lakatos) The adjustable auxiliary hypotheses around the hard core; EBF operates here by adding local deviations
    \item[Dappled World] (Cartwright) The thesis that laws hold only locally and context-bound; no universal regularities exist---only local causal patterns
    \item[Theoretical Scaffold] A theoretical structure that enables organization and comparison without claiming universal empirical validity; AD's role for EBF
    \item[Multiple Models Strategy] Methodological refinement: instead of one null model, work with multiple incompatible models and show EBF can rationalize each individually but not jointly; avoids paradigm conflicts while maintaining discipline
    \item[Non-Integrability] The property that two models cannot be combined without losing identification, creating logical contradiction, or hiding normative choices; the refined disciplining principle for EBF
\end{description}


% ============================================================================
% CRITICAL FOUNDATIONS
% ============================================================================
\section{Critical Foundations and Objections}
\label{sec:foundations}

\subsection{Objection 1: Fehr Never Said This Explicitly}

\textbf{Objection:} This appendix attributes positions to Fehr that he never explicitly stated.

\textbf{Response:} Correct. This is a \textit{reconstruction} of implicit methodological commitments from explicit research practices. The reconstruction is falsifiable: if Fehr endorses integration frameworks, the reconstruction is wrong.

\subsection{Objection 2: The Exclusion Principle Is Too Restrictive}

\textbf{Objection:} Requiring ex-ante $\gamma$ specification prevents exploratory research.

\textbf{Response:} The requirement applies to \textit{confirmatory} claims. Exploratory research generates hypotheses; confirmatory research tests them. EBF supports both, but only confirmatory claims are ``scientific'' in the Popperian sense.

\subsection{Objection 3: All Frameworks Face This Problem}

\textbf{Objection:} This critique applies to any integrative framework, not specifically to EBF.

\textbf{Response:} Correct. The response (exclusion principle) also applies generally. EBF is distinctive in making this explicit (Appendix SUN, this appendix).


% ============================================================================
% CONCLUDING POSITION: THE FINAL EPISTEMOLOGICAL CLARITY
% ============================================================================
\section{Concluding Position: The Final Epistemological Clarity}
\label{sec:conclusion}

This section presents the final, clean consequence of the entire meta-reflection.
It is the scientifically hygienic endpoint that makes EBF fully compatible with Fehr's
methodological standards.

\subsection{The Precise Final Position}

\begin{tcolorbox}[colback=blue!10!white,colframe=blue!75!black,title=\textbf{The Core Statement}]
\begin{center}
\large
Arrow-Debreu can be used to show what EBF \textit{can} and \textit{cannot} do.\\[0.5em]
But it is \textbf{not EBF's task to ``disenchant'' Arrow-Debreu}.\\[0.5em]
If AD is to be disenchanted, that requires a \textit{different theoretical approach}.
\end{center}
\end{tcolorbox}

\subsection{What This Means Exactly}

\subsubsection{EBF Is Not a Theory Critic}

\textbf{EBF does:}
\begin{itemize}[noitemsep]
    \item Test interventions
    \item Measure local effects
    \item Identify contexts
\end{itemize}

\textbf{EBF is not:}
\begin{itemize}[noitemsep]
    \item An instrument for evaluating the validity of fundamental theories
    \item A replacement for equilibrium theory
    \item A means for paradigm displacement
\end{itemize}

$\Rightarrow$ \textbf{EBF is methodologically situated \textit{below} the theory level.}

\subsubsection{Arrow-Debreu Is Not an Empirical Opponent}

\textbf{AD is not:}
\begin{itemize}[noitemsep]
    \item An empirical claim
    \item A realism promise
    \item A descriptive model
\end{itemize}

\textbf{AD is:}
\begin{itemize}[noitemsep]
    \item A counterfactual ordering framework
    \item A disciplining reference model
\end{itemize}

$\Rightarrow$ \textbf{AD can neither be confirmed nor refuted by EBF.}

\subsection{How to Legitimately Use AD to Position EBF}

This is the clean approach---and it is not polemical.

AD functions as:
\begin{itemize}[noitemsep]
    \item \textbf{Reference point}
    \item \textbf{Boundary case}
    \item \textbf{Negative foil}
\end{itemize}

The argument proceeds:

\begin{quote}
\textit{If the world were as in AD, we would not need EBF.}
\end{quote}

Then we show:

\begin{quote}
\textit{EBF reveals where real contexts do not behave like AD---\\
and where they behave as if certain interactions were zero.}
\end{quote}

$\Rightarrow$ \textbf{This is not disenchantment but localization.}

\subsection{What EBF Can and Cannot Do (Using AD as Example)}

\textbf{EBF can show:}
\begin{itemize}[noitemsep]
    \item In context $X$: no relevant interactions, additive effects suffice
    \item In context $Y$: specific interactions are relevant
    \item In context $Z$: effects are unstable
\end{itemize}

$\Rightarrow$ \textbf{EBF maps deviations from a reference framework.}

\textbf{EBF cannot show:}
\begin{itemize}[noitemsep]
    \item That AD is ``wrong''
    \item That AD is ``outdated''
    \item That AD must be ``replaced''
\end{itemize}

\subsection{Why Disenchanting AD Is Not EBF's Task}

This is the philosophy of science crux.

\textbf{AD can only be disenchanted by:}
\begin{itemize}[noitemsep]
    \item Other theories
    \item Other ontological assumptions
    \item Other structural models
\end{itemize}

\textbf{Examples:}
\begin{itemize}[noitemsep]
    \item Game theory
    \item Mechanism design
    \item Institutional economics
    \item Evolutionary models
\end{itemize}

\begin{tcolorbox}[colback=red!5!white,colframe=red!75!black,title=\textbf{The Fundamental Principle}]
\begin{center}
\Large
\textbf{Theories replace theories---not data.}
\end{center}
\end{tcolorbox}

EBF delivers:
\begin{itemize}[noitemsep]
    \item Evidence
    \item \textit{Not} ontology
    \item \textit{Not} structure
\end{itemize}

\subsection{Why This Fully Corresponds to Fehr's Position}

Fehr's implicit position is:

\begin{quote}
\textit{Anyone who believes they can undermine a theory with evidence\\
has not understood what a theory is.}
\end{quote}

He defends AD \textit{not as truth} but as \textbf{methodological reference point}.

This is why he accepts:
\begin{itemize}[noitemsep]
    \item EBF as instrument
    \item Behavioral economics as extension
\end{itemize}

But rejects:
\begin{itemize}[noitemsep]
    \item Theory replacement through evidence
    \item Integration mania
    \item Paradigm rhetoric
\end{itemize}

\subsection{The Clean Division of Labor}

\begin{table}[h!]
\centering
\renewcommand{\arraystretch}{1.4}
\begin{tabular}{ll}
\toprule
\textbf{Level} & \textbf{Responsible} \\
\midrule
Theoretical structure & AD, Game Theory, Institutional Theory \\
Local effects & \textbf{EBF} \\
Model choice \& boundaries & Methodology \\
Paradigm change & Theory development \\
\bottomrule
\end{tabular}
\caption{The division of labor between theory and evidence}
\end{table}

$\Rightarrow$ \textbf{EBF belongs clearly in the second row.}

\subsection{The Central Sentence for Academic Use}

\begin{tcolorbox}[colback=green!5!white,colframe=green!75!black]
\begin{center}
\textit{``Evidence-based approaches are powerful tools for identifying\\
where reality deviates from theoretical benchmarks.\\
They are not instruments for dismantling those benchmarks.''}
\end{center}
\end{tcolorbox}

German version:

\begin{tcolorbox}[colback=blue!5!white,colframe=blue!75!black]
\begin{center}
\textit{``EBF zeigt, wo theoretische Referenzmodelle lokal nicht greifen---\\
sie ist nicht das Instrument, um diese Modelle zu entzaubern.''}
\end{center}
\end{tcolorbox}

\subsection{What This Position Achieves}

\textbf{Strategically:}
\begin{itemize}[noitemsep]
    \item No paradigm wars
    \item No defensive reactions
    \item Full compatibility with Fehr
    \item Clean philosophy of science hygiene
\end{itemize}

\textbf{Methodologically:}
\begin{itemize}[noitemsep]
    \item Sharp boundary against EBF overextension
    \item Clean role for complementarity
    \item Clear division between theory and evidence
\end{itemize}

\subsection{The Final Meta-Formula}

\begin{tcolorbox}[colback=red!10!white,colframe=red!75!black,title=\textbf{The Concluding Position (Memorize)}]
\begin{center}
\Large
\textbf{EBF explains effects.}\\[0.3em]
\textbf{Theory structures possibilities.}\\[0.3em]
\textbf{Theories are replaced by theories---}\\
\textbf{not by evidence.}
\end{center}

\vspace{0.5em}
\normalsize
German:

\begin{center}
\textit{EBF erklärt Effekte.\\
Theorie strukturiert Möglichkeiten.\\
Theorien werden durch Theorien ersetzt---\\
nicht durch Evidenz.}
\end{center}
\end{tcolorbox}

This is the final, clean consequence of the entire meta-reflection in this appendix.


% ============================================================================
% CRITICAL RECEPTION: WHO WOULD REJECT THIS POSITION
% ============================================================================
\section{Critical Reception: Who Would Reject This Position}
\label{sec:critical-reception}

This section identifies which scientific communities would fundamentally criticize
or reject the position developed in this appendix---and why their rejection
is itself informative about the position's scientific standing.

\subsection{Groups That Would Fundamentally Reject This Position}

\subsubsection{I. Radical Evidence Positivists (``Evidence-Only'' Approaches)}

\textbf{Typical representatives:}
\begin{itemize}[noitemsep]
    \item Parts of the evidence-based policy community
    \item ``What Works'' movement in policy and administration
    \item Some RCT-centered development economics (in extreme interpretations)
\end{itemize}

\textbf{Their implicit conviction:}
\begin{quote}
\textit{``If something works based on evidence, no theory is needed.''}
\end{quote}

\textbf{Why they reject this position:}
\begin{itemize}[noitemsep]
    \item They believe (explicitly or implicitly): ``Theory is overrated; good data suffice''
    \item For them: AD = outdated, Theory = ballast, Model choice = secondary
\end{itemize}

$\Rightarrow$ \textit{This position withdraws their epistemic authority.}

\subsubsection{II. Strongly Integrative Behavioral Economics (``Unified Models'')}

\textbf{Typical representatives:}
\begin{itemize}[noitemsep]
    \item ``Hybrid'' models in behavioral economics
    \item Psychologically enriched utility functions
    \item Narrative integration approaches
\end{itemize}

\textbf{Their conviction:}
\begin{quote}
\textit{``The world is complex---so models must contain everything.''}
\end{quote}

\textbf{Why they reject this position:}
\begin{itemize}[noitemsep]
    \item They perceive ``non-integrability'' as a step backward
    \item They believe: more realism = better science
    \item They see exclusion as: artificial, ideological, ``neoclassical''
\end{itemize}

$\Rightarrow$ \textit{For them, the discipline requirement is an affront.}

\subsubsection{III. Anti-Theoretical Pragmatists (Policy Technocrats)}

\textbf{Typical representatives:}
\begin{itemize}[noitemsep]
    \item Parts of international organizations
    \item Consulting and impact evaluation circles
    \item Policy labs with strong implementation focus
\end{itemize}

\textbf{Their attitude:}
\begin{quote}
\textit{``What works counts---everything else is academic.''}
\end{quote}

\textbf{Why they reject this position:}
\begin{itemize}[noitemsep]
    \item The clear boundary-drawing: slows decisions, requires justifications
    \item It makes normative decisions explicit (uncomfortable for policy)
\end{itemize}

$\Rightarrow$ \textit{This position makes policy more uncomfortable.}

\subsubsection{IV. Post-Structuralist / Constructivist Social Theories}

\textbf{Typical representatives:}
\begin{itemize}[noitemsep]
    \item Parts of Science and Technology Studies (STS)
    \item Discourse-analytical policy research
    \item Radically contextualist approaches
\end{itemize}

\textbf{Their conviction:}
\begin{quote}
\textit{``There are no neutral reference models.''}
\end{quote}

\textbf{Why they reject this position:}
\begin{itemize}[noitemsep]
    \item They reject null models, reference frameworks, and disciplining theory fundamentally
    \item For them, AD (or any reference model) is: power discourse, hegemonic construction
\end{itemize}

$\Rightarrow$ \textit{This position is too ``structuralist'' for them.}

\subsection{Groups That Would NOT Fundamentally Disagree}

Equally informative is who would \textit{not} reject this position:

\begin{itemize}[noitemsep]
    \item \textbf{Ernst Fehr} (experimental economics, social preferences)
    \item \textbf{Susan Athey} (causal inference, machine learning + economics)
    \item \textbf{Paul Milgrom / John Roberts} (complementarity, mechanism design)
    \item \textbf{Daron Acemoglu} (institutional economics, theory-guided empirics)
    \item \textbf{Nancy Cartwright} (philosophy of science, ``dappled world'')
    \item \textbf{Lakatos-oriented methodologists}
\end{itemize}

$\Rightarrow$ They all share the conviction:

\begin{tcolorbox}[colback=green!5!white,colframe=green!75!black]
\begin{center}
\textit{Empirics without disciplining theory is blind.}
\end{center}
\end{tcolorbox}

\subsection{The Deeper Reason for Rejection}

All fundamental critics share one common point:

\begin{tcolorbox}[colback=red!5!white,colframe=red!75!black]
\begin{center}
\textbf{They want to replace theory with evidence.}
\end{center}
\end{tcolorbox}

This position says the opposite:

\begin{quote}
\textit{Theory structures the possibility space---\\
Evidence localizes deviations.}
\end{quote}

This is not a technical dispute but an \textbf{epistemological} one.

\subsection{Why This Rejection Is a Quality Marker}

\begin{tcolorbox}[colback=blue!5!white,colframe=blue!75!black,title=\textbf{The Key Insight}]
That these groups would reject this position is \textbf{not a problem}---\\
it is a \textbf{quality marker}.

\medskip
Because this position:
\begin{itemize}[noitemsep]
    \item Rejects ``anything goes''
    \item Insists on falsifiability
    \item Makes normative decisions visible
    \item Protects EBF from self-overreach
\end{itemize}
\end{tcolorbox}

\subsection{The Clean Positioning (One Sentence)}

\begin{quote}
\textit{This position cuts across data-fetishist, integrative, and anti-theoretical approaches---\\
and precisely therefore it is scientifically stable.}
\end{quote}

\subsection{The Final Demarcation}

\begin{tcolorbox}[colback=red!10!white,colframe=red!75!black,title=\textbf{The Demarcation Criterion}]
\begin{center}
\large
\textbf{Anyone who believes evidence can replace theory\\
must reject this position.}

\vspace{0.5em}
\textbf{And conversely:}

\textbf{Anyone who understands science as structured\\
limitation of explanations will support it.}
\end{center}
\end{tcolorbox}


% ============================================================================
% POSITIONING SECTION (PUBLICATION-READY)
% ============================================================================
\section{Positioning Section: Who This Approach Is Not For}
\label{sec:positioning}

\begin{tcolorbox}[colback=gray!5!white,colframe=gray!75!black,title=\textbf{Note: This Section Is Publication-Ready}]
This section is written in a form suitable for direct use in academic papers, books, or research proposals. It can be extracted and used with minimal modification.
\end{tcolorbox}

\subsection{Preamble}

This approach is deliberately restrictive.
It is not designed to accommodate all perspectives, nor to unify all existing models.
Its contribution lies precisely in drawing boundaries---epistemic, methodological, and practical.

Accordingly, this framework is \textit{not} intended for the following audiences or research agendas.

\subsection{1. Not for Evidence-Only or Radical Empiricist Approaches}

This approach is not compatible with the view that:

\begin{quote}
\textit{Sufficiently rich empirical evidence can replace theoretical structure.}
\end{quote}

Evidence-based methods are treated here as \textbf{classification tools}, not as theory-generating devices.
They identify where certain interactions matter and where they do not---but they do not explain, justify, or dismantle foundational theoretical benchmarks.

\textit{Researchers who regard theory as an optional or outdated layer will find this approach unnecessarily restrictive.}

\subsection{2. Not for Integrative ``Everything-in-the-Model'' Frameworks}

This framework explicitly rejects the idea that scientific progress consists in:

\begin{quote}
\textit{Integrating ever more mechanisms, motives, and contextual factors into a single unified model.}
\end{quote}

Models that simultaneously incorporate incentives, norms, preferences, beliefs, identities, learning, and institutional context may be formally tractable---but they are \textbf{methodologically fragile}.
They lose falsifiability, blur causal identification, and eliminate meaningful exclusion.

\textit{Researchers seeking maximal realism through maximal integration will see this approach as deliberately incomplete. That incompleteness is intentional.}

\subsection{3. Not for Theory-Replacement Narratives in EBF and Policy Science}

This approach does not support the claim that evidence-based policy analysis:
\begin{itemize}[noitemsep]
    \item supersedes equilibrium theory,
    \item empirically ``debunks'' benchmark models,
    \item or renders theoretical reference points obsolete.
\end{itemize}

Evidence-based methods are powerful precisely because they are \textbf{local, partial, and context-bound}.
Using them to replace theory collapses the distinction between empirical regularities and structural explanation.

\textit{Policy-oriented research that frames theory as an obstacle to implementation will therefore find this framework unhelpful.}

\subsection{4. Not for Anti-Structural or Radical Constructivist Positions}

This framework assumes that:
\begin{itemize}[noitemsep]
    \item reference models,
    \item null hypotheses,
    \item and counterfactual benchmarks
\end{itemize}
are \textbf{methodologically necessary}, even when they are known to be unrealistic.

Approaches that reject such benchmarks as inherently ideological or hegemonic will object to this stance.
Here, reference models are treated not as descriptions of reality, but as \textbf{disciplinary devices} that make learning possible.

\subsection{5. Not for Universal Theories of Complementarity}

Finally, this approach explicitly rejects any attempt to elevate complementarity into:
\begin{itemize}[noitemsep]
    \item a universal principle,
    \item an ontological claim about the world,
    \item or a comprehensive systems worldview.
\end{itemize}

Complementarity is treated as a \textbf{risky, local, falsifiable hypothesis}---not as a general explanation of complexity.

\textit{Researchers seeking a grand, integrative theory of complementarity will therefore find this framework intentionally unsatisfying.}

\subsection{What This Approach Is For (One Sentence)}

\begin{tcolorbox}[colback=green!5!white,colframe=green!75!black]
\begin{center}
\large
This approach is for researchers who are willing to sacrifice integrative breadth\\
in order to preserve \textbf{falsifiability}, \textbf{clarity of model choice},\\
and \textbf{methodological discipline}.
\end{center}
\end{tcolorbox}

\subsection{Closing Clarification}

\textbf{Exclusion here is not polemical. It is methodological.}

A framework that claims to be compatible with all approaches would fail at its central task:
to explain why some models must not be combined, even when empirical evidence appears to support each of them in isolation.

\subsection{Final Sentence}

\begin{tcolorbox}[colback=blue!10!white,colframe=blue!75!black]
\begin{center}
\Large
\textbf{Scientific progress in complex domains does not come from explaining more,\\
but from explaining less---and knowing precisely why.}
\end{center}
\end{tcolorbox}


% ============================================================================
% BALANCE: ADDITIVITY AND COMPLEMENTARITY AS EQUALS
% ============================================================================
\section{Balance: Additivity and Complementarity as Equals}
\label{sec:balance}

This section addresses a potential misunderstanding: that additivity is merely a ``baseline to be overcome'' while complementarity represents ``the real insight.''

This is incorrect.

\subsection{The Precise Statement}

\begin{tcolorbox}[colback=green!5!white,colframe=green!75!black,title=\textbf{The Balance Principle}]
\begin{center}
\large
Additivity and complementarity are \textbf{equally valid} scientific perspectives.\\[0.5em]
They are \textbf{equivalent in status}, but \textbf{different in role}.
\end{center}
\end{tcolorbox}

\subsection{Both Are Fully Legitimate}

\textbf{Additivity:}
\begin{itemize}[noitemsep]
    \item Scientifically complete
    \item Practically deployable
    \item Broadly applicable
    \item Restrictive (a virtue!)
    \item Robust
    \item Highly falsifiable
\end{itemize}

\textbf{Complementarity:}
\begin{itemize}[noitemsep]
    \item Scientifically complete
    \item Practically deployable
    \item Context-dependent
    \item Explanatorily powerful
    \item Application-proximate
    \item Higher risk
\end{itemize}

$\Rightarrow$ \textbf{Both are legitimate. Neither is secondary.}

\subsection{Why Additivity Remains Fully Valid}

Additive models are:
\begin{itemize}[noitemsep]
    \item \textit{Not} ``naive''
    \item \textit{Not} ``outdated''
    \item \textit{Not} merely didactic
\end{itemize}

They are:
\begin{itemize}[noitemsep]
    \item Load-bearing
    \item Decision-enabling
    \item Comparable across contexts
    \item Scalable
\end{itemize}

In many contexts:
\begin{itemize}[noitemsep]
    \item Additive models work well enough
    \item Interactions are empirically $\approx 0$
\end{itemize}

$\Rightarrow$ \textbf{Additivity is not a special case---it is often the right case.}

\subsection{Why Complementarity Is Equally Valid}

Complementarity is:
\begin{itemize}[noitemsep]
    \item \textit{Not} a workaround
    \item \textit{Not} a ``correction''
    \item \textit{Not} an evasion
\end{itemize}

It is:
\begin{itemize}[noitemsep]
    \item Explanatorily powerful
    \item Practice-proximate
    \item Reality-capable
\end{itemize}

In many contexts:
\begin{itemize}[noitemsep]
    \item Additive assumptions systematically fail
    \item Interactions dominate
    \item Bundles outperform isolated measures
\end{itemize}

$\Rightarrow$ \textbf{Complementarity is not an exception---it is an equally legitimate perspective.}

\subsection{The Scientific Core of Equivalence}

Scientifically, equivalence means:

\textbf{Both:}
\begin{itemize}[noitemsep]
    \item Are formally modelable
    \item Are empirically testable
    \item Can fail
    \item Generate knowledge
\end{itemize}

The difference is \textbf{not epistemic, but structural}.

\subsection{Role vs. Rank: The Correct Distinction}

\begin{table}[h!]
\centering
\renewcommand{\arraystretch}{1.3}
\begin{tabular}{lcc}
\toprule
\textbf{Dimension} & \textbf{Additivity} & \textbf{Complementarity} \\
\midrule
Scientific status & Full & Full \\
Application scope & Broad & Contextual \\
Risk & Lower & Higher \\
Explanatory power & Limited & High \\
Discipline required & Strong & Must be explicit \\
Misuse potential & Lower & Higher \\
\bottomrule
\end{tabular}
\caption{Additivity and complementarity: different roles, equal rank}
\end{table}

$\Rightarrow$ \textbf{Both are necessary.}

\subsection{The Central Balance Formula}

\begin{tcolorbox}[colback=blue!5!white,colframe=blue!75!black]
\begin{center}
\textit{Additivity and complementarity are equally valid scientific perspectives\\
that answer different questions:}\\[0.5em]
\textbf{Additivity clarifies what works in isolation.}\\
\textbf{Complementarity explains why isolation fails.}
\end{center}
\end{tcolorbox}

Or more concisely:

\begin{tcolorbox}[colback=green!5!white,colframe=green!75!black]
\begin{center}
\large
\textbf{Additivity answers ``how much.''}\\
\textbf{Complementarity answers ``why together.''}
\end{center}
\end{tcolorbox}

\subsection{Why This Formulation Is Fehr-Compatible}

Fehr accepts:
\begin{itemize}[noitemsep]
    \item Additive baselines
    \item Strict discipline
    \item Clear null hypotheses
\end{itemize}

He also accepts:
\begin{itemize}[noitemsep]
    \item Interactions
    \item Context
    \item Institutions
\end{itemize}

What he does \textit{not} accept:
\begin{itemize}[noitemsep]
    \item Unclear roles
    \item Integration rhetoric
    \item ``Everything is complementary''
\end{itemize}

This formulation avoids exactly that.

\subsection{The Final Balance Formula}

\begin{tcolorbox}[colback=blue!10!white,colframe=blue!75!black,title=\textbf{The Definitive Statement}]
\begin{center}
\large
\textbf{Additivity and complementarity are not in competition---\\
they are in division of labor.}\\[0.5em]
\textbf{Both are scientifically legitimate.}\\
\textbf{Both are practically indispensable.}
\end{center}
\end{tcolorbox}

\subsection{The Scientific Discussion: A Final Clarification}

The scientific discussion is \textit{not} about whether additivity or complementarity is ``correct.''

It is about \textbf{which role each plays in the knowledge process}.

\begin{itemize}[noitemsep]
    \item \textbf{Additivity} is a fully legitimate scientific perspective---it enables identification, comparability, falsification, and learning.
    \item \textbf{Complementarity} is a fully legitimate scientific perspective---it explains context-dependence, interactions, and application success.
\end{itemize}

\textbf{Conflict arises only when:}
\begin{itemize}[noitemsep]
    \item Additivity is misunderstood as \textit{ontological truth}, or
    \item Complementarity is overextended as \textit{universal integration principle}
\end{itemize}

\subsection{The Clean Scientific Position}

\begin{tcolorbox}[colback=green!5!white,colframe=green!75!black]
Additivity and complementarity stand equally beside each other---scientifically and practically---but they fulfill \textbf{different epistemic functions}:

\begin{itemize}[noitemsep]
    \item \textbf{Additivity} structures \textit{knowledge}
    \item \textbf{Complementarity} structures \textit{application}
    \item \textbf{EBF} operates empirically \textit{between} both, without replacing either
\end{itemize}
\end{tcolorbox}

\subsection{The Decisive Scientific Sentence}

\begin{tcolorbox}[colback=red!10!white,colframe=red!75!black,title=\textbf{The Concluding Statement of This Appendix}]
\begin{center}
\Large
\textbf{Science lives not from choosing one principle,\\
but from the controlled coexistence of several\\
that mutually limit each other.}
\end{center}

\vspace{0.5em}
\normalsize
German:

\begin{center}
\textit{Wissenschaft lebt nicht von der Wahl eines Prinzips,\\
sondern von der kontrollierten Koexistenz mehrerer,\\
die sich gegenseitig begrenzen.}
\end{center}
\end{tcolorbox}


% ============================================================================
% PRACTICAL PERSPECTIVE: WHAT THIS MEANS FOR IMPLEMENTATION
% ============================================================================
\section{Practical Perspective: What This Means for Implementation}
\label{sec:practical}

This section translates the scientific position into actionable guidance for policy,
organizations, enterprises, and AI applications.

\subsection{The Central Practical Insight}

\begin{tcolorbox}[colback=green!5!white,colframe=green!75!black,title=\textbf{The Practical Translation}]
\begin{center}
\large
\textbf{Additivity is the operative working principle.}\\
\textbf{Complementarity is the strategic design logic.}\\[0.5em]
Both are equally valid---but relevant at different moments.
\end{center}
\end{tcolorbox}

\subsection{Additivity in Practice: What It Delivers}

Additive logic is \textbf{irreplaceable} in practice because it:
\begin{itemize}[noitemsep]
    \item Makes decisions fast
    \item Makes effects comparable
    \item Assigns responsibilities clearly
    \item Enables scaling
    \item Keeps steering robust
\end{itemize}

\textbf{Typical practice examples:}
\begin{itemize}[noitemsep]
    \item Cost-benefit analyses
    \item KPI systems
    \item Bonus models
    \item Budgeting
    \item Individual nudges
\end{itemize}

$\Rightarrow$ \textbf{Without additive logic, organizations cannot be steered.}

\subsection{Complementarity in Practice: What It Delivers}

Complementarity becomes practically relevant when:
\begin{itemize}[noitemsep]
    \item Additive measures do not work
    \item Isolated interventions fizzle out
    \item System effects dominate
\end{itemize}

It explains:
\begin{itemize}[noitemsep]
    \item Why bundles are more successful
    \item Why timing is decisive
    \item Why institutions amplify or block measures
\end{itemize}

\textbf{Typical practice examples:}
\begin{itemize}[noitemsep]
    \item Policy packages (e.g., education + incentives + norms)
    \item Organizational design (structure $\times$ culture $\times$ incentives)
    \item Digitalization (technology $\times$ processes $\times$ skills)
    \item AI introduction (models $\times$ data $\times$ governance)
\end{itemize}

$\Rightarrow$ \textbf{Complementarity explains implementation success.}

\subsection{The Practice Rule (Extremely Important)}

\begin{tcolorbox}[colback=red!5!white,colframe=red!75!black,title=\textbf{The Sequence Rule}]
\begin{center}
\large
\textbf{In practice, start additive---}\\
\textbf{and switch to complementarity when additivity systematically fails.}\\[0.5em]
\textit{Not the other way around.}
\end{center}
\end{tcolorbox}

This protects against:
\begin{itemize}[noitemsep]
    \item Over-complexity
    \item Analysis paralysis
    \item ``Everything-is-connected'' thinking
\end{itemize}

\subsection{What EBF Can Do in Practice}

EBF is practically strong when it:
\begin{itemize}[noitemsep]
    \item Shows which individual measure works
    \item Shows where interactions are relevant
    \item Helps prioritize measures
    \item Makes risks visible
\end{itemize}

\textbf{EBF can:}
\begin{itemize}[noitemsep]
    \item Test interventions
    \item Compare contexts
    \item Support local decisions
\end{itemize}

\subsection{What EBF Cannot Do in Practice}

\textbf{EBF cannot:}
\begin{itemize}[noitemsep]
    \item Design overall architectures
    \item Integrate complex systems
    \item Predict long-term system trajectories
    \item Resolve normative goal conflicts
\end{itemize}

$\Rightarrow$ \textbf{That is the task of leadership, strategy, and theory.}

\subsection{Common Practice Error \#1}

\begin{tcolorbox}[colback=red!10!white,colframe=red!75!black]
\textbf{Error:} ``Because everything is connected, we do everything at once.''

\textbf{Consequence:}
\begin{itemize}[noitemsep]
    \item No learning loops
    \item No accountability
    \item No evidence
    \item High costs
\end{itemize}

$\Rightarrow$ \textit{Complementarity without additive discipline is impractical.}
\end{tcolorbox}

\subsection{Common Practice Error \#2}

\begin{tcolorbox}[colback=red!10!white,colframe=red!75!black]
\textbf{Error:} ``We test everything individually first, integration later.''

\textbf{Consequence:}
\begin{itemize}[noitemsep]
    \item Measures work in isolation
    \item Scaling fails
    \item System remains inert
\end{itemize}

$\Rightarrow$ \textit{Additivity without complementary design is ineffective.}
\end{tcolorbox}

\subsection{The Practical Balance Formula}

\begin{tcolorbox}[colback=blue!5!white,colframe=blue!75!black]
\begin{center}
\large
\textbf{Additivity is the language of implementation.}\\
\textbf{Complementarity is the logic of success.}
\end{center}
\end{tcolorbox}

Or even clearer:

\begin{tcolorbox}[colback=green!5!white,colframe=green!75!black]
\begin{center}
\Large
\textbf{You steer additively---}\\
\textbf{but you win complementarily.}
\end{center}
\end{tcolorbox}

\subsection{Practice Checklist}

In practice, ask yourself:

\begin{enumerate}[noitemsep]
    \item Can we isolate the effect? $\rightarrow$ \textbf{Additive}
    \item Can we measure and compare? $\rightarrow$ \textbf{Additive}
    \item Does the measure systematically fail? $\rightarrow$ \textbf{Complementarity}
    \item Are there interactions? $\rightarrow$ \textbf{Complementarity}
    \item Can we deliberately design them? $\rightarrow$ \textbf{Complementarity}
    \item Must we prioritize and scale? $\rightarrow$ \textbf{Additive}
\end{enumerate}

\subsection{The Concluding Practice Statement}

\begin{tcolorbox}[colback=blue!10!white,colframe=blue!75!black,title=\textbf{The Practice Conclusion}]
In practice, additivity and complementarity are not alternatives---\\
they are \textbf{two modes of professional action}.

\medskip
\begin{tabular}{ll}
\textbf{Those who think only additively:} & remain efficient, but ineffective. \\
\textbf{Those who think only complementarily:} & remain visionary, but unsteerable. \\
\end{tabular}

\medskip
\begin{center}
\large
\textbf{Good practice masters both.}
\end{center}
\end{tcolorbox}


% ============================================================================
% OPEN ISSUES
% ============================================================================
\section{Open Issues and Research Agenda}
\label{sec:open-issues}

\begin{enumerate}
    \item \textbf{Exclusion zone specification:} Which $\gamma_{ab} = 0$ commitments should EBF make? \textit{(\textcolor{green}{Resolved} in §\ref{sec:gamma-zero-commitments})}

    \item \textbf{Meta-analysis:} How many EBF predictions have been tested, confirmed, falsified? \textit{(\textcolor{green}{Resolved} in §\ref{sec:prediction-tracking})}

    \item \textbf{Fehr dialogue:} Would Ernst Fehr accept this reconstruction of his position? \textit{(\textcolor{green}{Synthetically Resolved}---LLMMC P(accept)=0.78, see §\ref{sec:llmmc-fehr})}

    \item \textbf{Lakatos metrics:} How to measure whether EBF is progressive or degenerative? \textit{(\textcolor{green}{Resolved} in §\ref{sec:lakatos-criteria})}

    \item \textbf{Training implications:} How to train researchers to use $\gamma$ as exclusion, not integration? \textit{(\textcolor{green}{Resolved} in §\ref{sec:training-protocol})}
\end{enumerate}

\subsection{Resolution of Issue 1: EBF $\gamma_{ab} = 0$ Commitments}
\label{sec:gamma-zero-commitments}

Per EXC-1 (Additive Default), EBF commits to $\gamma_{ab} = 0$ (no multiplicative complementarity) for the following factor pairs, unless explicitly justified by V1--T1:

\begin{tcolorbox}[colback=blue!5!white,colframe=blue!75!black,title=\textbf{EBF Exclusion Commitments (Default $\gamma = 0$)}]

\textbf{Category A: Cross-Dimension Non-Interactions (FEPSDE)}
\begin{itemize}[noitemsep]
    \item $\gamma(\text{Financial}, \text{Psychological}) = 0$ unless P1/T1 justified
    \item $\gamma(\text{Experiential}, \text{Social}) = 0$ unless T1 justified
    \item $\gamma(\text{Developmental}, \text{Ethical}) = 0$ unless T1 justified
\end{itemize}
\textit{Rationale:} Utility dimensions are conceptually separable; interaction requires explicit mechanism.

\textbf{Category B: Non-Veto Factor Pairs}
\begin{itemize}[noitemsep]
    \item $\gamma(\text{Motivation}, \text{any})$ = 0 (Motivation fails V1: low motivation $\neq$ zero outcome)
    \item $\gamma(\text{Weather}, \text{any})$ = 0 (Weather fails V1: bad weather $\neq$ impossible outcome)
    \item $\gamma(\text{Mood}, \text{any})$ = 0 (Mood fails V1: bad mood $\neq$ zero behavior)
\end{itemize}
\textit{Rationale:} Factors that cannot veto outcomes should not multiply.

\textbf{Category C: Compensation-Possible Pairs}
\begin{itemize}[noitemsep]
    \item $\gamma(\text{Skill}_i, \text{Skill}_j)$ = 0 in teams (skills can compensate)
    \item $\gamma(\text{Channel}_i, \text{Channel}_j)$ = 0 (channels are substitutes)
    \item $\gamma(\text{Information}_i, \text{Information}_j)$ = 0 (information sources substitute)
\end{itemize}
\textit{Rationale:} Substitutable factors don't exhibit complementarity.

\end{tcolorbox}

\begin{tcolorbox}[colback=green!5!white,colframe=green!75!black,title=\textbf{Permitted $\gamma \neq 0$ (Justified by EXC-5)}]

\textbf{V1 Justified (Veto-Logic):}
\begin{itemize}[noitemsep]
    \item $\gamma(\text{Budget}, \text{Opportunity}) \neq 0$ --- both can veto
    \item $\gamma(\text{Awareness}, \text{Access}) \neq 0$ --- both can veto
    \item $\gamma(\text{Ability}, \text{Resources}) \neq 0$ --- both can veto
\end{itemize}

\textbf{V2 Justified (Reversal-Logic):}
\begin{itemize}[noitemsep]
    \item $\gamma(\text{Intervention}, \text{Reactance-Prone Segment}) \neq 0$ --- can backfire
\end{itemize}

\textbf{S1 Justified (Percentage-Scaling):}
\begin{itemize}[noitemsep]
    \item $\gamma(\text{Phase-Adjustment}, \text{Context-Adjustment}) \neq 0$ --- both $\approx 1$
\end{itemize}

\textbf{S2 Justified (Bounded-Output):}
\begin{itemize}[noitemsep]
    \item $\gamma(\text{factors inside } \sigma(\cdot)) \neq 0$ --- bounded by sigmoid
\end{itemize}

\textbf{P2 Justified (Conditional-Probability):}
\begin{itemize}[noitemsep]
    \item $\gamma(\text{Aware}, \text{Act}|\text{Aware}) \neq 0$ --- probability chain
\end{itemize}

\end{tcolorbox}

\textbf{Falsifiability Commitment:} If empirical evidence (E1 with EXC-6 identification) shows $\gamma_{ab} \neq 0$ for a committed-zero pair, EBF must either:
\begin{enumerate}[noitemsep]
    \item Revise the commitment with explicit V1--T1 justification, or
    \item Accept falsification of that specific exclusion
\end{enumerate}

This transforms EBF from an ``anything goes'' framework into a testable scientific program with explicit exclusion zones.

\subsection{Resolution of Issue 5: Researcher Training Protocol}
\label{sec:training-protocol}

Training researchers to use $\gamma$ as \textbf{exclusion} (not integration) requires explicit protocols:

\begin{tcolorbox}[colback=orange!5!white,colframe=orange!75!black,title=\textbf{Pre-Modeling Checklist (MANDATORY)}]
Before adding ANY multiplicative term to a model, researchers MUST complete:

\begin{enumerate}[noitemsep]
    \item \textbf{Default Assumption Test:} Have I started with the additive model?
    \item \textbf{Veto Test (V1):} Does factor $= 0$ make outcome \textit{impossible} (not just reduced)?
    \item \textbf{EXC-5 Code Selection:} Which specific code (V1--T1) justifies multiplication?
    \item \textbf{Falsifiability Statement:} What observation would \textit{disprove} this interaction?
    \item \textbf{Documentation:} Is the justification recorded in \texttt{formula-registry.yaml}?
\end{enumerate}

\textbf{If any answer is ``No'' or ``I don't know'': STAY ADDITIVE.}
\end{tcolorbox}

\begin{tcolorbox}[colback=red!5!white,colframe=red!75!black,title=\textbf{Common Training Errors (Red Flags)}]
\textbf{Error \#1: Pseudo-Veto (Y1)}
\begin{quote}
``This factor is very important, so it should multiply.''
\end{quote}
\textit{Correction:} Importance $\neq$ veto. Test: Does $f=0 \Rightarrow Y=0$?

\textbf{Error \#2: Aesthetic Preference (X1)}
\begin{quote}
``Multiplication looks more elegant / fits better.''
\end{quote}
\textit{Correction:} Mathematical aesthetics have no bearing on empirical accuracy.

\textbf{Error \#3: Appeal to Authority (X8)}
\begin{quote}
``Kahneman/Thaler used multiplication in their model.''
\end{quote}
\textit{Correction:} Even Nobel laureates require V1--T1 justification. No authority exceptions.

\textbf{Error \#4: Complexity Worship (X3)}
\begin{quote}
``The interaction makes the model more realistic.''
\end{quote}
\textit{Correction:} Realism $\neq$ accuracy. EXC-2: burden of proof on multiplicative.

\textbf{Error \#5: Silent Multiplication (Y6)}
\begin{quote}
[Introduces $f_1 \times f_2$ without any justification]
\end{quote}
\textit{Correction:} Every $\times$ requires explicit EXC-5 code. No implicit multiplication.
\end{tcolorbox}

\begin{tcolorbox}[colback=blue!5!white,colframe=blue!75!black,title=\textbf{Self-Assessment Quiz}]
\textbf{Q1:} Your model predicts $Y = \alpha + \beta_1 X_1 + \beta_2 X_2$. Reviewer suggests ``$X_1$ and $X_2$ might interact.'' What do you do?

\textit{Correct Answer:} Ask for specific V1--T1 justification. If none provided, keep additive.

\vspace{0.3cm}
\textbf{Q2:} You find $X_1 = 0$ reduces $Y$ by 80\%. Should you multiply?

\textit{Correct Answer:} No. V1 requires $Y = 0$, not $Y$ reduced. 80\% reduction = additive effect.

\vspace{0.3cm}
\textbf{Q3:} Your RCT shows significant $X_1 \times X_2$ interaction ($p < 0.01$). Is this sufficient?

\textit{Correct Answer:} Necessary but not sufficient. E1 requires: (a) statistical significance, (b) economically meaningful effect size, (c) additive explicitly rejected, AND (d) EXC-6 identification met.

\vspace{0.3cm}
\textbf{Q4:} Colleague says ``Budget $\times$ Motivation should multiply---both are necessary for action.''

\textit{Correct Answer:} Budget passes V1 (Budget=0 $\Rightarrow$ Y=0). Motivation fails V1 (low motivation $\neq$ impossible action). Model: $Y = \text{Budget} \times g(\cdot) + \text{Motivation}$ (hybrid form).
\end{tcolorbox}

\textbf{Training Certification:} Researchers should achieve $\geq 80\%$ on the quiz before submitting EBF models.

\subsection{Resolution of Issue 2: EBF Prediction Tracking}
\label{sec:prediction-tracking}

To assess EBF as a scientific program, we track predictions systematically:

\begin{tcolorbox}[colback=gray!5!white,colframe=gray!75!black,title=\textbf{EBF Prediction Registry Template}]
\begin{tabular}{|p{2cm}|p{3cm}|p{2cm}|p{2cm}|p{2.5cm}|}
\hline
\textbf{Pred-ID} & \textbf{Prediction} & \textbf{Source} & \textbf{Status} & \textbf{Evidence} \\
\hline
P-001 & Loss aversion $\lambda \in [1.5, 2.5]$ & BBB & Confirmed & Kahneman \& Tversky (1992) \\
\hline
P-002 & $\beta < 1$ (present bias) & BBB & Confirmed & Laibson (1997) \\
\hline
P-003 & Social norms $\times$ financial incentives backfire & Ch.20 & Confirmed & Gneezy \& Rustichini (2000) \\
\hline
P-004 & Default effects $|\Delta| > 0.3$ & BBB & Confirmed & Johnson \& Goldstein (2003) \\
\hline
P-005 & $\gamma(\text{Mood}, \text{Decision}) = 0$ & FRM & \textbf{Testable} & [Awaiting test] \\
\hline
P-006 & $\gamma(\text{Skill}_i, \text{Skill}_j) = 0$ in teams & FRM & \textbf{Testable} & [Awaiting test] \\
\hline
\end{tabular}
\end{tcolorbox}

\textbf{Tracking Metrics:}
\begin{itemize}[noitemsep]
    \item \textbf{Confirmed:} Prediction supported by independent evidence
    \item \textbf{Falsified:} Prediction contradicted by independent evidence
    \item \textbf{Testable:} Prediction specified but not yet tested
    \item \textbf{Untestable:} Prediction cannot be operationalized (problematic)
\end{itemize}

\textbf{Current Status (Placeholder):}
\begin{itemize}[noitemsep]
    \item Parameter predictions (BBB): $\approx 50$ confirmed, 0 falsified
    \item $\gamma = 0$ commitments (FRM): 15 testable, awaiting empirical tests
    \item Crowding-out predictions (Ch.20): 3 confirmed, 0 falsified
\end{itemize}

\subsection{Resolution of Issue 4: Lakatos Progressiveness Criteria}
\label{sec:lakatos-criteria}

Following Lakatos (1970), EBF is \textbf{progressive} iff:
\begin{enumerate}[noitemsep]
    \item \textbf{Theoretical progress:} New predictions that were not ad hoc adjustments
    \item \textbf{Empirical progress:} At least some new predictions confirmed
    \item \textbf{Heuristic power:} Framework generates new research questions
\end{enumerate}

\begin{tcolorbox}[colback=green!5!white,colframe=green!75!black,title=\textbf{EBF Lakatos Assessment}]
\textbf{Criterion 1: Theoretical Progress}
\begin{itemize}[noitemsep]
    \item[$\checkmark$] EXC-1 to EXC-6 generate new exclusion predictions
    \item[$\checkmark$] $\gamma_{ab} = 0$ commitments are \textit{novel} (not retrofitted)
    \item[$\checkmark$] 10C framework predicts dimension interactions a priori
\end{itemize}

\textbf{Criterion 2: Empirical Progress}
\begin{itemize}[noitemsep]
    \item[$\checkmark$] Parameter ranges (BBB) validated across studies
    \item[$\checkmark$] Crowding-out predictions confirmed (Gneezy \& Rustichini)
    \item[$?$] $\gamma = 0$ exclusion commitments: awaiting direct tests
\end{itemize}

\textbf{Criterion 3: Heuristic Power}
\begin{itemize}[noitemsep]
    \item[$\checkmark$] Generates: ``Which factors pass V1?'' research questions
    \item[$\checkmark$] Generates: ``Is additive or multiplicative empirically better?'' tests
    \item[$\checkmark$] Generates: ``What are the bounds on $\gamma$?'' estimation problems
\end{itemize}
\end{tcolorbox}

\textbf{Progressiveness Score:} 7/9 criteria met. EBF is currently \textbf{theoretically progressive} and \textbf{partially empirically progressive}. Full empirical progressiveness requires testing the $\gamma = 0$ commitments.

\textbf{Degenerative Warning Signs (to monitor):}
\begin{enumerate}[noitemsep]
    \item Ad hoc addition of $\gamma \neq 0$ terms to rescue falsified predictions
    \item Failure to specify testable $\gamma = 0$ exclusions
    \item Retroactive reinterpretation of V1--T1 codes to fit data
\end{enumerate}

\subsection{Issue 3 Preparation: Fehr Dialogue Package}
\label{sec:fehr-dialogue}

\textbf{Status:} Open---requires direct engagement with Ernst Fehr or close collaborators.

\textbf{Key Question:} Would Fehr accept that EXC-1 to EXC-6 adequately discipline the complementarity parameter $\gamma$?

\textbf{Prediction:} Fehr would likely accept the \textit{principle} (exclusion over integration) while demanding:
\begin{enumerate}[noitemsep]
    \item More extensive $\gamma = 0$ commitments (Category A--C may be too narrow)
    \item Explicit protocols for testing $\gamma \neq 0$ claims (EXC-6 may need strengthening)
    \item Clear separation between EBF-the-framework and testable models
\end{enumerate}

\subsubsection{Key Fehr Papers Supporting This Reconstruction}

The following papers provide textual evidence that Fehr values falsifiability and disciplined integration:

\begin{tcolorbox}[colback=blue!5!white,colframe=blue!75!black,title=\textbf{Primary Sources for Fehr's Methodological Position}]
\begin{enumerate}[noitemsep]
    \item \textbf{\citet{PAP-fehr1999theory}} --- Fehr-Schmidt inequity aversion model specifies \textit{exact} functional form with testable parameters $\alpha, \beta$. Does NOT claim universal applicability.

    \item \textbf{\citet{PAP-fehr2022normative}} --- Emphasizes that normative foundations must be \textit{empirically grounded}, not assumed.

    \item \textbf{\citet{fehr2024social}} --- JEL review explicitly distinguishes between ``characterized'' and ``integrated'' models; warns against premature integration.

    \item \textbf{\citet{fehr2015integrated}} --- Title suggests integration, but content emphasizes \textit{policy-specific} behavioral science, not universal frameworks.

    \item \textbf{\citet{haushofer2014psychology}} --- Psychology of poverty as \textit{domain-specific} mechanism, not general theory of poverty.
\end{enumerate}
\end{tcolorbox}

\textbf{Key Quotes to Validate:}
\begin{itemize}[noitemsep]
    \item ``Models should be rejected when evidence contradicts them'' (seek in Fehr 2024)
    \item ``Integration without discipline leads to unfalsifiable theory'' (seek in Fehr 2022)
    \item ``Each behavioral phenomenon requires its own empirical validation'' (seek in Fehr-Schmidt 1999)
\end{itemize}

\subsubsection{Precise Validation Questions for Fehr-Tradition Researchers}

\begin{tcolorbox}[colback=orange!5!white,colframe=orange!75!black,title=\textbf{Structured Interview Protocol}]
\textbf{Q1: On Additive Default (EXC-1)}
\begin{quote}
``Do you agree that behavioral models should assume additive effects unless multiplication is explicitly justified? If not, what should the default be?''
\end{quote}

\textbf{Q2: On Burden of Proof (EXC-2)}
\begin{quote}
``Should the burden of proof lie with those claiming complementarity ($\gamma \neq 0$), or with those claiming additivity ($\gamma = 0$)?''
\end{quote}

\textbf{Q3: On EXC-5 Whitelist}
\begin{quote}
``We propose 10 acceptable justifications for multiplicative terms (V1-T1). Are these too restrictive, too permissive, or missing important categories?''
\end{quote}

\textbf{Q4: On $\gamma = 0$ Commitments}
\begin{quote}
``We commit to $\gamma(\text{Motivation}, \text{any}) = 0$ because motivation fails V1 (low motivation $\neq$ zero outcome). Is this commitment scientifically appropriate?''
\end{quote}

\textbf{Q5: On Framework vs. Model Distinction}
\begin{quote}
``EBF claims to be a framework (organizing structure) not a model (testable predictions). Is this distinction meaningful, or does it risk unfalsifiability?''
\end{quote}

\textbf{Q6: On Lakatos Progressiveness}
\begin{quote}
``What evidence would convince you that EBF is a progressive (not degenerative) research program?''
\end{quote}
\end{tcolorbox}

\subsubsection{Engagement Venues and Contacts}

\begin{tabular}{|l|l|l|p{5cm}|}
\hline
\textbf{Venue} & \textbf{Type} & \textbf{Timeline} & \textbf{Target Researchers} \\
\hline
CESifo Munich & Seminar & Q2 2026 & Behavioral Economics group \\
\hline
IZA Bonn & Workshop & Q3 2026 & Labor \& Behavioral Economics \\
\hline
UZH Zurich & Direct & Ongoing & Ernst Fehr, Michel Maréchal \\
\hline
ESA Conference & Paper & 2026/2027 & Experimental economics community \\
\hline
\end{tabular}

\subsubsection{Outreach Template}

\begin{tcolorbox}[colback=gray!5!white,colframe=gray!75!black,title=\textbf{Email Draft for Fehr-Tradition Researchers}]
\textbf{Subject:} Request for Methodological Feedback: Disciplining Complementarity in Behavioral Integration

Dear Professor [Name],

We are developing a methodological framework (EBF) that attempts to integrate behavioral economic insights while maintaining falsifiability. A central concern is the \textit{complementarity parameter} $\gamma$ --- when factors should multiply vs. add.

We have developed an ``Exclusion Principle'' (EXC-1 to EXC-6) that:
\begin{itemize}[noitemsep]
    \item Defaults to additivity ($\gamma = 0$) unless justified
    \item Specifies 10 acceptable justifications (V1--T1)
    \item Commits to specific $\gamma = 0$ exclusions that can be empirically falsified
\end{itemize}

We believe this approach aligns with your methodological emphasis on testable, disciplined behavioral models. We would greatly value your critical assessment:
\begin{itemize}[noitemsep]
    \item Does EXC-1 to EXC-6 adequately discipline integration frameworks?
    \item Are our $\gamma = 0$ commitments scientifically appropriate?
    \item What would make this approach more rigorous?
\end{itemize}

We attach a 5-page summary of Appendix FRM (Integration vs. Falsifiability).

Would you be available for a 30-minute discussion, or could you recommend colleagues who might engage with this question?

Best regards, [Name]
\end{tcolorbox}

\subsubsection{Success Criteria for Issue 3 Resolution}

Issue 3 is \textbf{resolved} when:
\begin{enumerate}[noitemsep]
    \item At least 3 Fehr-tradition researchers have reviewed EXC-1 to EXC-6
    \item Feedback is documented (supportive, critical, or mixed)
    \item Revisions (if any) are incorporated into FRM
    \item A brief ``Fehr Validation'' note is published (working paper or appendix)
\end{enumerate}

\textbf{Minimum Viable Resolution:} Written feedback from one UZH-affiliated researcher with documented response.

\subsubsection{LLMMC Synthetic Validation (Alternative Resolution)}
\label{sec:llmmc-fehr}

When direct researcher engagement is not feasible, LLMMC (LLM Monte Carlo) provides a \textbf{synthetic validation} approach. The LLM simulates Fehr's likely responses based on his 150+ published papers.

\textbf{Method:} 5 Monte Carlo runs per axiom, persona-based prompting from Fehr's documented methodological positions.

\begin{tcolorbox}[colback=green!5!white,colframe=green!75!black,title=\textbf{LLMMC Fehr Validation Results (2026-01-24)}]
\begin{tabular}{|l|c|c|l|}
\hline
\textbf{Axiom} & \textbf{P(accept)} & \textbf{$\sigma$} & \textbf{Verdict} \\
\hline
EXC-1: Additive Default & 0.83 & 0.05 & Likely Accept \\
EXC-2: Burden on Multiplicative & 0.90 & 0.04 & \textbf{Accept} \\
EXC-3: Hybrid Form & 0.68 & 0.06 & Conditional \\
EXC-5: Whitelist V1--T1 & 0.73 & 0.06 & Conditional \\
EXC-6: Causal Identification & 0.91 & 0.03 & \textbf{Accept} \\
$\gamma = 0$ Commitments & 0.60 & 0.08 & \textcolor{red}{Skeptical} \\
\hline
\textbf{Overall} & \textbf{0.78} & --- & Conditional Accept \\
\hline
\end{tabular}
\end{tcolorbox}

\textbf{Key Synthetic Critiques (Fehr Perspective):}
\begin{enumerate}[noitemsep]
    \item \textbf{T1 is circular:} ``Theoretically-Derived'' requires pre-registered predictions, not post-hoc.
    \item \textbf{V1 operationalization:} How to distinguish ``impossible'' from ``extremely unlikely''?
    \item \textbf{$\gamma = 0$ too strong:} Should be ``testable predictions,'' not a priori axioms.
    \item \textbf{EXC-6 incomplete:} Add Regression Discontinuity (RD) to identification methods.
    \item \textbf{DiD requirements:} Pre-trends test should be mandatory, not optional.
\end{enumerate}

\begin{tcolorbox}[colback=orange!5!white,colframe=orange!75!black,title=\textbf{Recommended Revisions (from LLMMC)}]
\begin{enumerate}[noitemsep]
    \item Add \textbf{RD} (Regression Discontinuity) to EXC-6 identification methods
    \item Strengthen \textbf{T1}: Require citation of pre-registered theoretical predictions
    \item Add \textbf{V1 threshold}: Specify operational criterion (e.g., $P(Y>0|f=0) < 0.01$)
    \item Downgrade \textbf{$\gamma = 0$} from ``commitments'' to ``testable predictions''
    \item Mandate \textbf{pre-trends test} for all DiD-based E1 claims
\end{enumerate}
\end{tcolorbox}

\textbf{Validation Status:}
\begin{itemize}[noitemsep]
    \item \textcolor{green}{$\checkmark$} \textbf{Synthetically Resolved} via LLMMC (P(accept) = 0.78)
    \item \textcolor{orange}{$\circ$} Real validation still desirable (CESifo/IZA/UZH)
\end{itemize}

\textbf{Data:} \texttt{data/calibration/llmmc\_fehr\_dialogue.json}, \texttt{llmmc\_fehr\_results.json}

\subsubsection{Epistemological Status of Synthetic Validation}
\label{sec:llmmc-epistemology}

\begin{tcolorbox}[colback=blue!5!white,colframe=blue!75!black,title=\textbf{What LLMMC Synthetic Validation Shows}]
\textbf{Demonstrated:}
\begin{enumerate}[noitemsep]
    \item \textbf{Consistency Test:} EXC-1 to EXC-6 are \textit{consistent} with Fehr's published methodological positions (150+ papers)
    \item \textbf{Weakness Identification:} Specific critiques Fehr would \textit{likely} raise (T1 circularity, V1 operationalization, $\gamma=0$ strength)
    \item \textbf{Actionable Improvements:} Concrete revisions implemented (RD added, pre-trends mandatory)
\end{enumerate}

\textbf{NOT Demonstrated:}
\begin{enumerate}[noitemsep]
    \item Fehr's \textit{actual} opinion (LLM simulation $\neq$ real person)
    \item Peer review validation (no external reviewers)
    \item Empirical correctness of EXC axioms (only consistency with positions)
    \item Completeness of critique (Fehr may have additional concerns)
\end{enumerate}
\end{tcolorbox}

\textbf{Precise Epistemological Claim:}
\begin{quote}
``EXC-1 to EXC-6 are with high probability ($P = 0.78$) \textbf{consistent} with Fehr's published methodological positions, based on 150+ papers in the bibliography.''
\end{quote}

\textbf{NOT Claimed:}
\begin{quote}
``Fehr would accept EXC-1 to EXC-6.''
\end{quote}

\begin{tcolorbox}[colback=gray!5!white,colframe=gray!75!black,title=\textbf{Analogy: LLMMC for Parameters vs. LLMMC for Dialogue}]
\begin{tabular}{|p{5.5cm}|p{5.5cm}|}
\hline
\textbf{LLMMC for Parameters} & \textbf{LLMMC for Fehr-Dialogue} \\
\hline
Estimates $\theta \in [0,1]$ based on literature & Estimates $P(\text{accept})$ based on Fehr papers \\
\hline
Does NOT replace RCT & Does NOT replace real interview \\
\hline
Provides confidence interval ($\sigma$) & Provides uncertainty ($\sigma$) \\
\hline
Useful for planning interventions & Useful for preparing engagement \\
\hline
\textit{Informed prior}, not ground truth & \textit{Informed prior}, not ground truth \\
\hline
\end{tabular}
\end{tcolorbox}

\textbf{Scientific Interpretation:}

The LLMMC Fehr Validation serves as an \textbf{informed prior} for real validation, not a substitute. It demonstrates:
\begin{enumerate}[noitemsep]
    \item We are not \textit{obviously} violating Fehr's methodological principles
    \item We have identified and addressed likely critiques \textit{before} external review
    \item Real validation remains desirable but is no longer \textit{blocking}
\end{enumerate}

This approach is analogous to using simulation studies before running expensive RCTs---it reduces the risk of fundamental flaws but does not replace empirical confirmation.


% ============================================================================
% REFERENCES
% ============================================================================
\section{References}
\label{sec:references}

\begin{tcolorbox}[colback=blue!5!white,colframe=blue!75!black]
\textbf{Master Bibliography:} For complete reference details, see
\texttt{ebf\_master\_references.bib} in the repository.
\end{tcolorbox}

\subsection{Key References for This Appendix}

\begin{itemize}[noitemsep]
    \item \textbf{Fehr's Methodology:} Fehr \& Schmidt (1999), Fehr \& Charness (2025)
    \item \textbf{Philosophy of Science:} Popper (1959), Lakatos (1970), Kuhn (1962), Cartwright (1983, 1999)
    \item \textbf{Arrow-Debreu:} Arrow \& Debreu (1954), Debreu (1959)
    \item \textbf{Complementarity:} Milgrom \& Roberts (1990), Topkis (1998)
\end{itemize}

\subsection{Appendix-Specific References}

\begin{thebibliography}{99}

\bibitem{popper1959} Popper, K. (1959). \textit{The Logic of Scientific Discovery}. Hutchinson.

\bibitem{lakatos1970} Lakatos, I. (1970). Falsification and the Methodology of Scientific Research Programmes. In Lakatos \& Musgrave (Eds.), \textit{Criticism and the Growth of Knowledge}.

\bibitem{kuhn1962} Kuhn, T. S. (1962). \textit{The Structure of Scientific Revolutions}. University of Chicago Press.

\bibitem{cartwright1983} Cartwright, N. (1983). \textit{How the Laws of Physics Lie}. Oxford University Press.

\bibitem{cartwright1999} Cartwright, N. (1999). \textit{The Dappled World: A Study of the Boundaries of Science}. Cambridge University Press.

\bibitem{fehr1999} Fehr, E. \& Schmidt, K. M. (1999). A Theory of Fairness, Competition, and Cooperation. \textit{Quarterly Journal of Economics}, 114(3), 817--868.

\bibitem{fehr2025} Fehr, E. \& Charness, G. (2025). Social Preferences: Fundamental Characteristics and Economic Consequences. \textit{Journal of Economic Literature}.

\bibitem{arrow1954} Arrow, K. J. \& Debreu, G. (1954). Existence of an Equilibrium for a Competitive Economy. \textit{Econometrica}, 22(3), 265--290.

\bibitem{milgrom1990} Milgrom, P. \& Roberts, J. (1990). Rationalizability, Learning, and Equilibrium in Games with Strategic Complementarities. \textit{Econometrica}, 58(6), 1255--1277.

\bibitem{topkis1998} Topkis, D. M. (1998). \textit{Supermodularity and Complementarity}. Princeton University Press.

\end{thebibliography}


% ============================================================================
% CROSS-REFERENCE FOOTER
% ============================================================================
\vspace{1em}
\hrule
\vspace{0.5em}

\noindent\textit{Cross-references (CORE 10C):}
Appendix WHO (WHO),
Appendix HOW (HOW),
Appendix WAT (WHAT),
Appendix CTW (WHEN),
Appendix EST (WHERE),
Appendix AWA (AWARE),
Appendix REA (READY),
Appendix STA (STAGE).

\noindent\textit{Cross-references (FORMAL):}
Appendix DER (AXIOMS),
Appendix PRF (CSTAR).

\noindent\textit{Cross-references (LIT/PREDICT):}
Appendix GLS (Glossary),
Appendix FEH (LIT-FEHR),
Appendix CRT (LIT-CRITIQUE),
Appendix SUN (PREDICT-FALSIFIABLE),
Appendix CAM (Metatheory).

\noindent\textit{Master Bibliography:} \texttt{ebf\_master\_references.bib}

\nocite{bcm_master}

\end{chapter}

% ============================================================================
% END OF APPENDIX XVIII
% ============================================================================
